\section{High-Speed Autonomous PIV Data Templates (Informative)}

This appendix illustrates the per-entry authentication key selector, per-card
certificate discovery, and explicit resumption of PIV Auto. The configuration,
discovery, resume, and TWIC sequences added in this revision are illustrative;
they have not been executed or used to regenerate cryptographic vectors.
The normative rules are in Subsections~\ref{sec:pivauto-key-discovery},
\ref{sec:pivauto-twic}, and~\ref{sec:pivauto-resume}.

The existing APDU fixtures and generated signing examples are retained separately.
Their prototype KDK input uses the SC1-style 48-byte session-key set and does not
establish Secure Channel~2.0 interoperability. SC2 remains a required upstream
dependency; its transport and final PIV Auto binding are outside these examples.

\subsection{Vector Sources}
\begin{SimpleTable}{PIV Auto Vector Sources}{3}
  {|L{0.24\SimpleUsableWidth}|L{0.22\SimpleUsableWidth}|L{0.54\SimpleUsableWidth}|}
  {\TableHeaderCell{Source} & \TableHeaderCell{Use} & \TableHeaderCell{Notes}}
  NIST SD 33 / IR 8347 & Live-card APDU fixtures & Cards 2, 3, 4, 5, and 16 provide VCI contactless enrolment, OPACITY, and key-challenge transcripts. \\\hline
  NIST SM/VCI captures & CVC and trust anchors & Provides validator-ready \ApduTemplate{7F21}, \ApduTemplate{5FC122}, and \ApduTemplate{7F50} objects. \\\hline
  GSA FICAM card-builder & Local keys and certificates & Used only as reusable local test material for source-backed simulation. \\\hline
  Generated YubiKey PIV & Missing algorithm profiles & Reusable RSA-1024, RSA-3072, and P-384 test keys are committed as public fixtures and loaded through \code{ykman piv}; RSA-1024 is informative legacy test material only. \\\hline
\end{SimpleTable}

\subsection{Key Selection and Certificate Discovery}

An Auto-enabled Configuration Entry ends with one Authentication Key Reference
byte, after the AID bytes. There is no configured algorithm byte. With Auto
disabled, the key-reference byte is absent and the previous entry encoding is
unchanged. For example, the default PIV application with contact and contactless
permitted is \HexBytes{01040103009E}: tag \code{0x01}, length \code{0x04},
Data Model \code{0x01}, Interface Mask \code{0x03}, AID Length \code{0x00},
and key \code{0x9E}. Changing only the last byte to \code{0x9A} selects PIV
Authentication without fixing the card's key type.

The logical configuration begins with three header bytes, two four-byte counters,
and 32 entropy bytes. The shortest Auto-enabled configuration is therefore
\code{0x31} bytes (49 decimal), including its one six-byte Configuration Entry.
Configuration Count counts Configuration Entry TLVs only.

\begin{SimpleTable}{Illustrative PIV Auto Configuration and Resume Bytes}{3}
  {|L{0.25\SimpleUsableWidth}|L{0.27\SimpleUsableWidth}|K{0.48\SimpleUsableWidth}|}
  {\TableHeaderCell{Message} & \TableHeaderCell{Field} & \TableHeaderCell{Bytes}}
  Card Authentication & Flags, Count, Global Mask, Series, Sequence & \HexBytes{010103070000002A000000} \\\hline
  Card Authentication & Supplemental Entropy & \HexBytes{0000000000000000000000000000000000000000000000000000000000000000} \\\hline
  Card Authentication & Configuration Entry & \HexBytes{01040103009E} \\\hline
  Contactless PIV Authentication & Flags, Count, Global Mask, Series, Sequence & \HexBytes{030102070000002A000000} \\\hline
  Contactless PIV Authentication & Supplemental Entropy & \HexBytes{0000000000000000000000000000000000000000000000000000000000000000} \\\hline
  Contactless PIV Authentication & Configuration Entry & \HexBytes{01040102009A} \\\hline
  Resume after pairing and PIN & Flags, Count, Global Mask, Series, Sequence & \HexBytes{030102070000002B000000} \\\hline
  Resume after pairing and PIN & Supplemental Entropy & \HexBytes{0000000000000000000000000000000000000000000000000000000000000000} \\\hline
  Resume after pairing and PIN & Unchanged Configuration Entry & \HexBytes{01040102009A} \\\hline
  Resume after pairing and PIN & Final Resume TLV & \HexBytes{0208070000002A000000} \\\hline
\end{SimpleTable}

Concatenate the rows belonging to one message. The two initial logical payloads
are \code{0x31} bytes; the resume payload is \code{0x3B} bytes. For a single
fragment, prefix the initial payload with \HexBytes{AE310000003100}, or the
resume payload with \HexBytes{AE3B0000003B00}. These prefixes contain CMND,
MpSizeTotal, MpOffset, and MpFragmentSize. OSDP packet headers and SC2 protection
are not shown. All-zero entropy is permitted but not recommended for deployment.

\begin{SimpleTable}{Illustrative Authentication Certificate Reads}{3}
  {|L{0.18\SimpleUsableWidth}|K{0.44\SimpleUsableWidth}|L{0.38\SimpleUsableWidth}|}
  {\TableHeaderCell{Configured Key} & \TableHeaderCell{Logical Card APDU} & \TableHeaderCell{Certificate Object}}
  PIV \code{0x9A} & \HexBytes{00CB3FFF055C035FC10500} & \code{0x5FC105}; contactless requires VCI. \\\hline
  PIV \code{0x9E} & \HexBytes{00CB3FFF055C035FC10100} & \code{0x5FC101}; no PIV PIN required. \\\hline
  NextGen TWIC \code{0x9E} & \HexBytes{00CB3FFF055C035FC101FF} & \code{0x5FC101} in the selected TWIC application. \\\hline
\end{SimpleTable}

These are logical APDUs; an established credential SM session protects applicable
PIV commands. The PD reassembles the returned certificate object, extracts tag
\code{0x70}, and interprets its one-byte CertInfo in tag \code{0x71}:
\code{0x00} identifies uncompressed DER and \code{0x01} identifies GZIP-compressed
DER. It decompresses when indicated before parsing SubjectPublicKeyInfo.
Missing or malformed certificate evidence ends the attempt with Result
\code{0x03}; a recognized but unsupported subject-key profile uses Result
\code{0x04}. Before an algorithm is determined, PIVSTATUS reports Algorithm
Identifier \code{0x00}; that failure value never selects a default card algorithm.

\begin{SimpleTable}{Mixed PIV Subject Keys under the Same Configuration}{4}
  {|L{0.27\SimpleUsableWidth}|L{0.23\SimpleUsableWidth}|K{0.18\SimpleUsableWidth}|L{0.32\SimpleUsableWidth}|}
  {\TableHeaderCell{Subject Public Key} & \TableHeaderCell{Configured Key} & \TableHeaderCell{Discovered Algorithm} & \TableHeaderCell{Challenge Profile}}
  RSA-2048 & \code{0x9A} or \code{0x9E} & \code{0x07} & 32-byte digest; SHA-256 DigestInfo. \\\hline
  RSA-3072 & \code{0x9A} or \code{0x9E} & \code{0x05} & 48-byte digest; SHA-384 DigestInfo. \\\hline
  Named curve P-256 & \code{0x9A} or \code{0x9E} & \code{0x11} & 32-byte prehashed ECDSA input. \\\hline
  Named curve P-384 & \code{0x9A} or \code{0x9E} & \code{0x14} & 48-byte prehashed ECDSA input. \\\hline
\end{SimpleTable}

The four subject-key types at each of the two references give eight PIV profiles.
Within either reference, successive RSA and ECC cards use identical configuration
bytes. The certificate changes the GENERAL AUTHENTICATE P1 byte, challenge
encoding, and reported algorithm. For example, an RSA CA can sign a P-384
subject certificate using SHA-256: the PD still discovers \code{0x14} and uses
the 48-byte P-384 profile. The issuer's signature algorithm does not select the
card's authentication algorithm. No alternate key is tried if discovery fails.

\subsection{Illustrative PIV Authentication Resume}

In the contactless \code{0x9A} example above, attempt \code{0x00000007}/
\code{0x0000002A} ends with Result \code{0x05}, Detail \code{0x01} because
pairing is required. The ACU receives the complete result, then sends
\texttt{osdp\_PIVXMITPAIRING}, for example
\HexBytes{AD3132333435363738}, and waits for confirmed VCI establishment.
It requests commanded PIN acquisition and verification with
\HexBytes{AA010000000000}: Mode \code{0x01}, default finite timeouts,
Auto Execute clear, and Cache Validity zero. The ACU waits for the successful
PIVSPER before sending the resume PIVMODE payload above.

The final \HexBytes{0208070000002A000000} identifies the completed prerequisite
failure, not the counters of the new attempt. The header retains Series
\code{0x00000007} and supplies its exact next unused Sequence
\code{0x0000002B}. Configuration Count remains one. The entropy may change;
the flags, masks, ordered entries, AIDs, and key references remain unchanged.
For this eligible resume, completed pairing and PIN verification remain bound
to the same presentation, application, interface, and security contexts.

The PD reads certificate \code{0x5FC105}, discovers its subject-key profile,
and performs one new key proof. Pairing or PIN completion alone does not restart
Auto. Credential removal, a context change, or an interrupted security session
invalidates resumption; an ordinary mode replacement still clears state.
Duplicate OSDP delivery does not repeat pairing, PIN verification, or the proof.

\subsection{Illustrative TWIC Auto Key Proof}

The TWIC Configuration Entry \HexBytes{010D020209A000000367200000019E} selects
Data Model \code{0x02}, contactless interface, the nine-byte TWIC AID prefix,
and asymmetric key \code{0x9E}. It does not configure an algorithm or a TPK.
The PD issues \HexBytes{00A4040009A0000003672000000100} and examines the full
AID returned in tag \code{0x4F}:
\begin{itemize}
  \item \textbf{Legacy:} \HexBytes{A000000367200000010101}. Read the TWIC
    CHUID with \HexBytes{00CB3FFF055C035FC102FF}, retaining its identity only
    for this attempt. Then select the PIV application with
    \HexBytes{00A404000BA00000030800001000010000} and retrieve certificate
    \code{0x5FC101}. Use that application's \code{0x9E} key while retaining
    the configured TWIC acceptance policy. A present, correctly sized nil GUID
    remains zero in the challenge and PIVSTATUS only when the valid nonzero
    FASC-N matches the Card Authentication certificate. Missing identity data
    does not receive the nil-GUID exception.
  \item \textbf{NextGen:} \HexBytes{A000000367200000010103}. Remain in the
    TWIC application, read CHUID \code{0x5FC102} and certificate
    \code{0x5FC101} using Le \code{0xFF}, and retain the populated GUID.
    Concatenate the initial data and subsequent GET RESPONSE data when status
    \code{0x61XX} indicates more bytes. Le \code{0x00} can return
    \code{0x9000} with no data and is not a successful certificate read.
\end{itemize}

The documented NextGen profile uses RSA-2048: discover \code{0x07} from the
certificate SPKI and use GENERAL AUTHENTICATE with P1 \code{0x07} and P2
\code{0x9E}. Neither workflow submits a PIV PIN or uses the TPK for this proof.
The Legacy application transition preserves the card-authentication purpose;
it does not authorize fallback from PIV Authentication \code{0x9A} to
Card Authentication \code{0x9E}. A TWIC-model \code{0x9A} entry is invalid.

These new TWIC examples contain no captured card responses or newly computed
signatures. Implementation fixtures must use public synthetic keys, certificates,
and invented identifiers, with their provenance kept separate from physical-card
captures. PIVSTATUS success reports the OSDP-defined key proof, not completion
of a TSA reader mode or an access decision. Certificate trust, identity binding,
expiration, and applicable cancellation checks remain required by deployment
policy. TWIC Part~3, Section~4.4.4 requires at least 127 bytes of unique challenge
data; these illustrative examples do not establish that fixed RSA padding and a
shorter KMAC digest satisfy that requirement. TSA reader qualification remains
a separate verification item.

\subsection{Supported Profiles}
\begin{SimpleTable}{PIV Auto Supported Profiles}{5}
  {|L{0.18\SimpleUsableWidth}|K{0.10\SimpleUsableWidth}|K{0.15\SimpleUsableWidth}|L{0.12\SimpleUsableWidth}|L{0.45\SimpleUsableWidth}|}
  {\TableHeaderCell{Profile} & \TableHeaderCell{Alg} & \TableHeaderCell{Key Ref} & \TableHeaderCell{PIV Auto} & \TableHeaderCell{Coverage}}
  RSA-1024 & \ApduHeader{06} & \ApduHeader{9A}, \ApduHeader{9E} & No & Informative legacy/non-current test-only material. \\\hline
  RSA-2048 & \ApduHeader{07} & \ApduHeader{9A}, \ApduHeader{9E} & Yes & NIST SD 33 cards 2, 3, 16; generated-card; GSA-backed simulation. \\\hline
  RSA-3072 & \ApduHeader{05} & \ApduHeader{9A}, \ApduHeader{9E} & Yes & Committed YubiKey material; requires YubiKey firmware with RSA-3072 support. \\\hline
  ECC P-256 & \ApduHeader{11} & \ApduHeader{9A}, \ApduHeader{9E} & Yes & NIST SD 33 cards 4 and 5; generated-card coverage. \\\hline
  ECC P-384 & \ApduHeader{14} & \ApduHeader{9A}, \ApduHeader{9E} & Yes & Committed YubiKey material. NIST card 5 also provides \ApduHeader{9C}/\ApduHeader{9D} APDU examples. \\\hline
\end{SimpleTable}

\subsection{Fixed-Width KMAC Challenge Template}
\begin{SimpleTable}{PIV Auto KMAC Challenge Template}{4}
  {|L{0.23\SimpleUsableWidth}|L{0.10\SimpleUsableWidth}|L{0.21\SimpleUsableWidth}|L{0.46\SimpleUsableWidth}|}
  {\TableHeaderCell{Field} & \TableHeaderCell{Size} & \TableHeaderCell{Encoding} & \TableHeaderCell{Notes}}
  Template label & 0x0E & ASCII plus version byte & \code{OSDP-PIV-AUTO} followed by \ApduChallenge{01}. \\\hline
  Series Counter & 0x04 & LSB first & ACU-provided series, for example \ApduChallenge{07000000}. \\\hline
  Sequence Counter & 0x04 & LSB first & PD-allocated attempt counter within that series, for example \ApduChallenge{2A000000}; see Subsection~\ref{sec:pivauto-counters}. \\\hline
  FASC-N & 0x19 & Raw bytes & Credential FASC-N from CHUID. \\\hline
  UUID & 0x10 & Raw bytes & Credential UUID; the recognized Legacy TWIC nil-GUID exception follows Subsection~\ref{sec:pivauto-twic}. Unavailable UUID otherwise causes failure. \\\hline
  PIV Key Reference & 0x01 & Raw byte & Challenged key reference, normally \ApduHeader{9A} or \ApduHeader{9E}. \\\hline
  PIV Algorithm & 0x01 & Raw byte & Discovered supported identifier: \ApduHeader{05}, \ApduHeader{07}, \ApduHeader{11}, or \ApduHeader{14}; \ApduHeader{06} remains informative legacy test material only. \\\hline
  Supplemental Entropy & 0x20 & Raw bytes & Fixed-width ACU entropy field; all zeroes are permitted but not recommended. \\\hline
\end{SimpleTable}

The fixed-width encoding and profile mapping are normative in
Tables~\ref{tab:pivauto-challenge} and~\ref{tab:pivauto-profiles}; this appendix
illustrates them. The retained prototype vectors derive their KDK as
\code{PIV-AUTO-KDK = KMAC256(key = S-ENC || S-MAC1 || S-MAC2, data = empty, mac\_len = 32, custom = "OSDP-PIV-AUTO-KDK-v1")}.
This is SC1-style prototype input, not evidence of an SC2 key binding.
The challenge stage is
\code{KMAC256(key = PIV-AUTO-KDK, data = fixed-width template, mac\_len = L, custom = "OSDP-PIV-AUTO-CHALLENGE-v1")}.
\code{L} is 32 bytes for SHA-256 profiles and 48 bytes for SHA-384 profiles.
The supplemental entropy field is always present and fixed at 32 bytes; all
zeroes are valid for test vectors but not recommended for deployments.

\subsection{Dynamic Authentication Templates}
\subsubsection{RSA Key Challenge Input}
\begin{SimpleTable}{RSA Dynamic Authentication Template}{4}
  {|K{0.16\SimpleUsableWidth}|L{0.16\SimpleUsableWidth}|L{0.26\SimpleUsableWidth}|L{0.42\SimpleUsableWidth}|}
  {\TableHeaderCell{Tag / Bytes} & \TableHeaderCell{Length} & \TableHeaderCell{Name} & \TableHeaderCell{Value}}
  \ApduTemplate{7C} & Variable & Dynamic Authentication Template & Outer template for the PIV GENERAL AUTHENTICATE command data. \\\hline
  \ApduResponse{82} & \ApduResponse{00} & Response request & Empty value requesting the card response in tag \ApduResponse{82}. \\\hline
  \ApduTemplate{81} & Modulus length & Challenge block & EMSA-PKCS1-v1\_5 encoded digest block with the same length as the RSA modulus. \\\hline
  \ApduChallenge{0001FF...FF00} & Variable & Padding prefix & EMSA-PKCS1-v1\_5 block prefix before the digest information. \\\hline
  \ApduChallenge{30...} & Variable & DigestInfo & ASN.1 DigestInfo for SHA-256 or SHA-384 followed by the KMAC-derived digest; SHA-256 uses prefix \ApduChallenge{3031300D060960864801650304020105000420} and SHA-384 uses prefix \ApduChallenge{3041300D060960864801650304020205000430}. \\\hline
\end{SimpleTable}

\begin{SimpleTable}{RSA Dynamic Authentication Examples}{6}
  {|L{0.16\SimpleUsableWidth}|K{0.11\SimpleUsableWidth}|K{0.11\SimpleUsableWidth}|L{0.18\SimpleUsableWidth}|L{0.20\SimpleUsableWidth}|L{0.24\SimpleUsableWidth}|}
  {\TableHeaderCell{Profile} & \TableHeaderCell{Alg} & \TableHeaderCell{Key Ref} & \TableHeaderCell{Tag 81} & \TableHeaderCell{Response 82} & \TableHeaderCell{Command APDU}}
  RSA-1024 & \ApduHeader{06} & \ApduHeader{9A}/\ApduHeader{9E} & 0x80 bytes & 0x80 bytes & Informative legacy/non-current test-only APDU. \\\hline
  RSA-2048 & \ApduHeader{07} & \ApduHeader{9A}/\ApduHeader{9E} & 0x100 bytes & 0x100 bytes & 0x112-byte extended APDU. \\\hline
  RSA-3072 & \ApduHeader{05} & \ApduHeader{9A}/\ApduHeader{9E} & 0x180 bytes & 0x180 bytes & 0x192-byte extended APDU. \\\hline
\end{SimpleTable}

\subsubsection{ECC Signature Input}
\begin{SimpleTable}{ECC Dynamic Authentication Template}{4}
  {|K{0.16\SimpleUsableWidth}|L{0.16\SimpleUsableWidth}|L{0.26\SimpleUsableWidth}|L{0.42\SimpleUsableWidth}|}
  {\TableHeaderCell{Tag / Bytes} & \TableHeaderCell{Length} & \TableHeaderCell{Name} & \TableHeaderCell{Value}}
  \ApduTemplate{7C} & Variable & Dynamic Authentication Template & Outer template for the PIV GENERAL AUTHENTICATE command data. \\\hline
  \ApduResponse{82} & \ApduResponse{00} & Response request & Empty value requesting the card response in tag \ApduResponse{82}. \\\hline
  \ApduTemplate{81} & Digest length & Challenge digest & KMAC-derived digest supplied directly to the ECC signing operation. \\\hline
  Digest bytes & 0x20 or 0x30 bytes & Digest value & SHA-256-width digest for P-256; SHA-384-width digest for P-384. \\\hline
\end{SimpleTable}

\begin{SimpleTable}{ECC Dynamic Authentication Examples}{6}
  {|L{0.16\SimpleUsableWidth}|K{0.11\SimpleUsableWidth}|K{0.11\SimpleUsableWidth}|L{0.18\SimpleUsableWidth}|L{0.20\SimpleUsableWidth}|L{0.24\SimpleUsableWidth}|}
  {\TableHeaderCell{Profile} & \TableHeaderCell{Alg} & \TableHeaderCell{Key Ref} & \TableHeaderCell{Tag 81} & \TableHeaderCell{Response 82} & \TableHeaderCell{Command APDU}}
  ECC P-256 & \ApduHeader{11} & \ApduHeader{9A}/\ApduHeader{9E} & 0x20 bytes & DER ECDSA signature, about 0x46-0x48 bytes & 0x2C-byte short APDU. \\\hline
  ECC P-384 & \ApduHeader{14} & \ApduHeader{9A}/\ApduHeader{9E} & 0x30 bytes & DER ECDSA signature, about 0x66-0x68 bytes & 0x3C-byte short APDU. \\\hline
\end{SimpleTable}

\subsubsection{Worked Dynamic Authentication Examples}
% Generated by scripts/render_piv_auto_appendix_tables.py; do not hand edit.
\TableNote{Worked examples are generated from \code{test-vectors/piv-auto-demo/generated/piv-auto-simulation-report.json}.}

\begin{SimpleTable}{PIV Auto Worked Example: PIV Authentication RSA-1024}{2}
  {|L{0.24\SimpleUsableWidth}|K{0.76\SimpleUsableWidth}|}
  {\TableHeaderCell{Field} & \TableHeaderCell{Value}}
  Profile & \code{9a-rsa1024} \\\hline
  Source & Committed YubiKey loader material in test-vectors/piv-auto-yubikey-material/9a-rsa1024. Informative legacy/non-current test-only coverage. \\\hline
  Algorithm ID & \code{0x}\ApduHeader{06} \\\hline
  Key Reference & \code{0x}\ApduHeader{9A} \\\hline
  KMAC Challenge Digest & \ApduChallenge{5B2FC08860E3999EEF8F6CF5D2B0C9AE05CBB28A5B423BFE4F646F6F47FCA4F2} \\\hline
  Tag 81 Challenge Input & \ApduChallenge{0001FFFFFFFFFFFFFFFFFFFFFFFFFFFFFFFFFFFFFFFFFFFFFFFFFFFFFFFFFFFFFFFFFFFFFFFFFFFFFFFFFFFFFFFFFFFFFFFFFFFFFFFFFFFFFFFFFFFFFFFFFFFFFFFFFFFFFFFFFFFFFFFFFFFF003031300D0609608648016503040201050004205B2FC08860E3999EEF8F6CF5D2B0C9AE05CBB28A5B423BFE4F646F6F47FCA4F2} \\\hline
  GENERAL AUTHENTICATE APDU & \ApduHeader{0087069A}\ApduLength{88}\ApduTemplate{7C8185}\ApduResponse{8200}\ApduTemplate{818180}\ApduChallenge{0001FFFFFFFFFFFFFFFFFFFFFFFFFFFFFFFFFFFFFFFFFFFFFFFFFFFFFFFFFFFFFFFFFFFFFFFFFFFFFFFFFFFFFFFFFFFFFFFFFFFFFFFFFFFFFFFFFFFFFFFFFFFFFFFFFFFFFFFFFFFFFFFFFFFF003031300D0609608648016503040201050004205B2FC08860E3999EEF8F6CF5D2B0C9AE05CBB28A5B423BFE4F646F6F47FCA4F2}\ApduLength{00} \\\hline
  Dynamic Authentication Response & \ApduTemplate{7C8183}\ApduResponse{828180}\ApduResponse{6FE396FF8D6F618EDFF35CC600ECE5DFECB74C66D6FE9035201271D2230484F989A70619293F8450AE67AF864C8C1F35AFBDC8D531D28AB6E3E8335588E123C224F6737807DCD3DCF78EDD6853FE216C2FEBF78DFCD7242650954781DEF3FBBC6C772E6CB65D1135051CD116FBAB39E3BDFBACA0A9949C6BE5189290BA83498A}\ApduStatus{9000} \\\hline
  Response Tag 0x82 Length & 0x0080 bytes \\\hline
\end{SimpleTable}

\begin{SimpleTable}{PIV Auto Worked Example: Card Authentication RSA-1024}{2}
  {|L{0.24\SimpleUsableWidth}|K{0.76\SimpleUsableWidth}|}
  {\TableHeaderCell{Field} & \TableHeaderCell{Value}}
  Profile & \code{9e-rsa1024} \\\hline
  Source & Committed YubiKey loader material in test-vectors/piv-auto-yubikey-material/9e-rsa1024. Informative legacy/non-current test-only coverage. \\\hline
  Algorithm ID & \code{0x}\ApduHeader{06} \\\hline
  Key Reference & \code{0x}\ApduHeader{9E} \\\hline
  KMAC Challenge Digest & \ApduChallenge{DB6851587148B203545E4306315583B0BB0CB88B8534200FDE187CF4F8F37D4F} \\\hline
  Tag 81 Challenge Input & \ApduChallenge{0001FFFFFFFFFFFFFFFFFFFFFFFFFFFFFFFFFFFFFFFFFFFFFFFFFFFFFFFFFFFFFFFFFFFFFFFFFFFFFFFFFFFFFFFFFFFFFFFFFFFFFFFFFFFFFFFFFFFFFFFFFFFFFFFFFFFFFFFFFFFFFFFFFFFF003031300D060960864801650304020105000420DB6851587148B203545E4306315583B0BB0CB88B8534200FDE187CF4F8F37D4F} \\\hline
  GENERAL AUTHENTICATE APDU & \ApduHeader{0087069E}\ApduLength{88}\ApduTemplate{7C8185}\ApduResponse{8200}\ApduTemplate{818180}\ApduChallenge{0001FFFFFFFFFFFFFFFFFFFFFFFFFFFFFFFFFFFFFFFFFFFFFFFFFFFFFFFFFFFFFFFFFFFFFFFFFFFFFFFFFFFFFFFFFFFFFFFFFFFFFFFFFFFFFFFFFFFFFFFFFFFFFFFFFFFFFFFFFFFFFFFFFFFF003031300D060960864801650304020105000420DB6851587148B203545E4306315583B0BB0CB88B8534200FDE187CF4F8F37D4F}\ApduLength{00} \\\hline
  Dynamic Authentication Response & \ApduTemplate{7C8183}\ApduResponse{828180}\ApduResponse{23F95C29066CB6513196C034958EF17FF08135027485AC6C91EE230226E7A81F446CB1CAD207C2B4D8F32C6BFD78824297270DAD5C66388208B868C790120CCC247E134C193AA7DE60CB6E66A5AC104867B0BEAB0468C8558B26E45FED226B4DA7F0EF9F8339A531CEDF354BEA70094F01E01267D0EEC359B4768FA360469C87}\ApduStatus{9000} \\\hline
  Response Tag 0x82 Length & 0x0080 bytes \\\hline
\end{SimpleTable}

\begin{SimpleTable}{PIV Auto Worked Example: PIV Authentication RSA-2048}{2}
  {|L{0.24\SimpleUsableWidth}|K{0.76\SimpleUsableWidth}|}
  {\TableHeaderCell{Field} & \TableHeaderCell{Value}}
  Profile & \code{9a-rsa2048} \\\hline
  Source & Committed RSA-2048 card-authentication key in test-vectors/piv-auto-demo/source-p12, originally imported from GSA FICAM card-builder. NIST SD 33 / IR 8347 and generated-card coverage. \\\hline
  Algorithm ID & \code{0x}\ApduHeader{07} \\\hline
  Key Reference & \code{0x}\ApduHeader{9A} \\\hline
  KMAC Challenge Digest & \ApduChallenge{E5D733C7FBBA8344701EE8909F1E6642E040CCF3265E0D5D28A26D73C344C19C} \\\hline
  Tag 81 Challenge Input & \ApduChallenge{0001FFFFFFFFFFFFFFFFFFFFFFFFFFFFFFFFFFFFFFFFFFFFFFFFFFFFFFFFFFFFFFFFFFFFFFFFFFFFFFFFFFFFFFFFFFFFFFFFFFFFFFFFFFFFFFFFFFFFFFFFFFFFFFFFFFFFFFFFFFFFFFFFFFFFFFFFFFFFFFFFFFFFFFFFFFFFFFFFFFFFFFFFFFFFFFFFFFFFFFFFFFFFFFFFFFFFFFFFFFFFFFFFFFFFFFFFFFFFFFFFFFFFFFFFFFFFFFFFFFFFFFFFFFFFFFFFFFFFFFFFFFFFFFFFFFFFFFFFFFFFFFFFFFFFFFFFFFFFFFFFFFFFFFFFFFFFFFFFFFFFFFFFFFFFFFFFFFFFFFFFFFFFFFFFFFFFFFFFFFFFFFFFFFFFFFFFFFFFFFFFFFFF003031300D060960864801650304020105000420E5D733C7FBBA8344701EE8909F1E6642E040CCF3265E0D5D28A26D73C344C19C} \\\hline
  GENERAL AUTHENTICATE APDU & \ApduHeader{0087079A}\ApduLength{00010A}\ApduTemplate{7C820106}\ApduResponse{8200}\ApduTemplate{81820100}\ApduChallenge{0001FFFFFFFFFFFFFFFFFFFFFFFFFFFFFFFFFFFFFFFFFFFFFFFFFFFFFFFFFFFFFFFFFFFFFFFFFFFFFFFFFFFFFFFFFFFFFFFFFFFFFFFFFFFFFFFFFFFFFFFFFFFFFFFFFFFFFFFFFFFFFFFFFFFFFFFFFFFFFFFFFFFFFFFFFFFFFFFFFFFFFFFFFFFFFFFFFFFFFFFFFFFFFFFFFFFFFFFFFFFFFFFFFFFFFFFFFFFFFFFFFFFFFFFFFFFFFFFFFFFFFFFFFFFFFFFFFFFFFFFFFFFFFFFFFFFFFFFFFFFFFFFFFFFFFFFFFFFFFFFFFFFFFFFFFFFFFFFFFFFFFFFFFFFFFFFFFFFFFFFFFFFFFFFFFFFFFFFFFFFFFFFFFFFFFFFFFFFFFFFFFFFF003031300D060960864801650304020105000420E5D733C7FBBA8344701EE8909F1E6642E040CCF3265E0D5D28A26D73C344C19C}\ApduLength{00} \\\hline
  Dynamic Authentication Response & \ApduTemplate{7C820104}\ApduResponse{82820100}\ApduResponse{515356AB0012D988701AC373B029D86540C4DF473B031EDB5ADFDB84BC214A45497F6CEE5FFB15ECB0CD89B9F351EF5006AB20105EF289EE041F427802867BA48253B6B90F0BD6F77637B14DD36B32BE89F8BD12F2D9F0E5B408C6C452959F233EDB9FA0A96B2808711FB1B671A683460F5FD62847AC18E5717B7D008EDE5FB2AEDFFF44EDD90B2E2DB595456BA7D42644C57C2267C4B9961424839E352FE64EB5F344CBB12BC630FFCDE943055A80375CF6AEA11BAF6AEC2916D5A1B6D769CC5F935519634C2504A2CA4916368DAAA887DC61A514AC630DB093CD2AC3608CECD355A63DCA02E35FBF2CAEDA34ED51C7F60C69EFE7CCCFDA9115A93950016990}\ApduStatus{9000} \\\hline
  Response Tag 0x82 Length & 0x0100 bytes \\\hline
\end{SimpleTable}

\begin{SimpleTable}{PIV Auto Worked Example: Card Authentication RSA-2048}{2}
  {|L{0.24\SimpleUsableWidth}|K{0.76\SimpleUsableWidth}|}
  {\TableHeaderCell{Field} & \TableHeaderCell{Value}}
  Profile & \code{9e-rsa2048} \\\hline
  Source & Committed RSA-2048 card-authentication key in test-vectors/piv-auto-demo/source-p12, originally imported from GSA FICAM card-builder. NIST SD 33 / IR 8347 and generated-card coverage. \\\hline
  Algorithm ID & \code{0x}\ApduHeader{07} \\\hline
  Key Reference & \code{0x}\ApduHeader{9E} \\\hline
  KMAC Challenge Digest & \ApduChallenge{0864C776F2374124D3E63F0B0B29FC1C5F0E8FF8BB1FA80E2723293B86A0158E} \\\hline
  Tag 81 Challenge Input & \ApduChallenge{0001FFFFFFFFFFFFFFFFFFFFFFFFFFFFFFFFFFFFFFFFFFFFFFFFFFFFFFFFFFFFFFFFFFFFFFFFFFFFFFFFFFFFFFFFFFFFFFFFFFFFFFFFFFFFFFFFFFFFFFFFFFFFFFFFFFFFFFFFFFFFFFFFFFFFFFFFFFFFFFFFFFFFFFFFFFFFFFFFFFFFFFFFFFFFFFFFFFFFFFFFFFFFFFFFFFFFFFFFFFFFFFFFFFFFFFFFFFFFFFFFFFFFFFFFFFFFFFFFFFFFFFFFFFFFFFFFFFFFFFFFFFFFFFFFFFFFFFFFFFFFFFFFFFFFFFFFFFFFFFFFFFFFFFFFFFFFFFFFFFFFFFFFFFFFFFFFFFFFFFFFFFFFFFFFFFFFFFFFFFFFFFFFFFFFFFFFFFFFFFFFFFFF003031300D0609608648016503040201050004200864C776F2374124D3E63F0B0B29FC1C5F0E8FF8BB1FA80E2723293B86A0158E} \\\hline
  GENERAL AUTHENTICATE APDU & \ApduHeader{0087079E}\ApduLength{00010A}\ApduTemplate{7C820106}\ApduResponse{8200}\ApduTemplate{81820100}\ApduChallenge{0001FFFFFFFFFFFFFFFFFFFFFFFFFFFFFFFFFFFFFFFFFFFFFFFFFFFFFFFFFFFFFFFFFFFFFFFFFFFFFFFFFFFFFFFFFFFFFFFFFFFFFFFFFFFFFFFFFFFFFFFFFFFFFFFFFFFFFFFFFFFFFFFFFFFFFFFFFFFFFFFFFFFFFFFFFFFFFFFFFFFFFFFFFFFFFFFFFFFFFFFFFFFFFFFFFFFFFFFFFFFFFFFFFFFFFFFFFFFFFFFFFFFFFFFFFFFFFFFFFFFFFFFFFFFFFFFFFFFFFFFFFFFFFFFFFFFFFFFFFFFFFFFFFFFFFFFFFFFFFFFFFFFFFFFFFFFFFFFFFFFFFFFFFFFFFFFFFFFFFFFFFFFFFFFFFFFFFFFFFFFFFFFFFFFFFFFFFFFFFFFFFFFF003031300D0609608648016503040201050004200864C776F2374124D3E63F0B0B29FC1C5F0E8FF8BB1FA80E2723293B86A0158E}\ApduLength{00} \\\hline
  Dynamic Authentication Response & \ApduTemplate{7C820104}\ApduResponse{82820100}\ApduResponse{9BB1AA3875DB4E41107DD91CEAAC36FE4CE5F969E1F7B9B199DFAB38CFA1EE36FBF197CD2777C4A6CE082623C9F36A9DC2514FE54DD12934AFCAE16DFE14FFCFC0F58E0A0F46567444968F1DE78B496B4222DB913A87B850139B1E2A730124C4ED169919F1EB5FA3E15617E89BB9CED8FA996A63E126C239EB6D142A5746FBEC3DFA3B287F944EE45586BA5AEB64087BDA0661DF7254C42D9768141F7243382D61CC43A8CDBE4ADA849E05F013D947DF7FA608C38ABA3841B038AA69D77461B00053F59C4D34906BFBFD470A277D0722AAECD83117DE906BFCD36A16EFB00B3EACD4414F48528CCA4BE57826B4312424A65DABAD6FA1ECA3DAF55F6E6E1022C8}\ApduStatus{9000} \\\hline
  Response Tag 0x82 Length & 0x0100 bytes \\\hline
\end{SimpleTable}

\begin{SimpleTable}{PIV Auto Worked Example: PIV Authentication RSA-3072}{2}
  {|L{0.24\SimpleUsableWidth}|K{0.76\SimpleUsableWidth}|}
  {\TableHeaderCell{Field} & \TableHeaderCell{Value}}
  Profile & \code{9a-rsa3072} \\\hline
  Source & Committed YubiKey loader material in test-vectors/piv-auto-yubikey-material/9a-rsa3072. Generated-card coverage. \\\hline
  Algorithm ID & \code{0x}\ApduHeader{05} \\\hline
  Key Reference & \code{0x}\ApduHeader{9A} \\\hline
  KMAC Challenge Digest & \ApduChallenge{1503FD1E89AA91A5B37E6872AC3846F654B7B1D099B2F3AEB597472DFE38E71442816BC43454622219F61A3C6DAFDB4F} \\\hline
  Tag 81 Challenge Input & \ApduChallenge{0001FFFFFFFFFFFFFFFFFFFFFFFFFFFFFFFFFFFFFFFFFFFFFFFFFFFFFFFFFFFFFFFFFFFFFFFFFFFFFFFFFFFFFFFFFFFFFFFFFFFFFFFFFFFFFFFFFFFFFFFFFFFFFFFFFFFFFFFFFFFFFFFFFFFFFFFFFFFFFFFFFFFFFFFFFFFFFFFFFFFFFFFFFFFFFFFFFFFFFFFFFFFFFFFFFFFFFFFFFFFFFFFFFFFFFFFFFFFFFFFFFFFFFFFFFFFFFFFFFFFFFFFFFFFFFFFFFFFFFFFFFFFFFFFFFFFFFFFFFFFFFFFFFFFFFFFFFFFFFFFFFFFFFFFFFFFFFFFFFFFFFFFFFFFFFFFFFFFFFFFFFFFFFFFFFFFFFFFFFFFFFFFFFFFFFFFFFFFFFFFFFFFFFFFFFFFFFFFFFFFFFFFFFFFFFFFFFFFFFFFFFFFFFFFFFFFFFFFFFFFFFFFFFFFFFFFFFFFFFFFFFFFFFFFFFFFFFFFFFFFFFFFFFFFFFFFFFFFFFFFFFFFFFFFFFFFFFFFFFFFFFFFFFFFFFFFFFFFFFFFFFFFFFFFFFFFFFFFFFFFFFFFFFFFFFFFFFFFFFFFFFFFFFFFFFFFFFFFFFFFFFFFFFFFF003041300D0609608648016503040202050004301503FD1E89AA91A5B37E6872AC3846F654B7B1D099B2F3AEB597472DFE38E71442816BC43454622219F61A3C6DAFDB4F} \\\hline
  GENERAL AUTHENTICATE APDU & \ApduHeader{0087059A}\ApduLength{00018A}\ApduTemplate{7C820186}\ApduResponse{8200}\ApduTemplate{81820180}\ApduChallenge{0001FFFFFFFFFFFFFFFFFFFFFFFFFFFFFFFFFFFFFFFFFFFFFFFFFFFFFFFFFFFFFFFFFFFFFFFFFFFFFFFFFFFFFFFFFFFFFFFFFFFFFFFFFFFFFFFFFFFFFFFFFFFFFFFFFFFFFFFFFFFFFFFFFFFFFFFFFFFFFFFFFFFFFFFFFFFFFFFFFFFFFFFFFFFFFFFFFFFFFFFFFFFFFFFFFFFFFFFFFFFFFFFFFFFFFFFFFFFFFFFFFFFFFFFFFFFFFFFFFFFFFFFFFFFFFFFFFFFFFFFFFFFFFFFFFFFFFFFFFFFFFFFFFFFFFFFFFFFFFFFFFFFFFFFFFFFFFFFFFFFFFFFFFFFFFFFFFFFFFFFFFFFFFFFFFFFFFFFFFFFFFFFFFFFFFFFFFFFFFFFFFFFFFFFFFFFFFFFFFFFFFFFFFFFFFFFFFFFFFFFFFFFFFFFFFFFFFFFFFFFFFFFFFFFFFFFFFFFFFFFFFFFFFFFFFFFFFFFFFFFFFFFFFFFFFFFFFFFFFFFFFFFFFFFFFFFFFFFFFFFFFFFFFFFFFFFFFFFFFFFFFFFFFFFFFFFFFFFFFFFFFFFFFFFFFFFFFFFFFFFFFFFFFFFFFFFFFFFFFFFFFFFFFFFF003041300D0609608648016503040202050004301503FD1E89AA91A5B37E6872AC3846F654B7B1D099B2F3AEB597472DFE38E71442816BC43454622219F61A3C6DAFDB4F}\ApduLength{00} \\\hline
  Dynamic Authentication Response & \ApduTemplate{7C820184}\ApduResponse{82820180}\ApduResponse{19C03ACCABE74C9259FA5A2267BC4F1401AB352F210BEFE0F74B8B74F25C82E0A83E1635C19C341CBD97D71D856BF132A2C3CDD223AC7D37451E9083A0C5FDB8013577741520A8A2375896E3839252F10DD8826A00A065F06154EE7A51703D65CC00AAA66E28BE5683085B4045882AA17B925457D3B6B0BC0C4581D10ED4653A1D4713E4FB5A4C72054789C9A4C1E6BC2B6B497228A9C765365040D00C139671175C262BA2E6D6064F1F84AAC0D5279B91F9FDE6D6B02B2423C7CC3FA7408E80A8040B92B819F80D6AD23B86C8F339D50F332802CBED0511215B03DF9496DD64F77ADBE4AE844F0D708F230E6F06911DADCA6888B9C80C94A16FCB0DE7467B70766668CDEA0053289295256F8BEDC693F058F6244BB6B315C4C9F07D89432A479AF72C666181C7B47CAA72C7ABE8B323658BD681FA10A1E88E2AFAB03BFDE9B571CA6FEDA05A24CC6436061D633691822FDEDD3DAE4E130A1CCE50E65FDFC5EA451B7AA8F83D4385D40E005E368D064E105D6F5325B36D955F3D88B5BEF84078}\ApduStatus{9000} \\\hline
  Response Tag 0x82 Length & 0x0180 bytes \\\hline
\end{SimpleTable}

\begin{SimpleTable}{PIV Auto Worked Example: Card Authentication RSA-3072}{2}
  {|L{0.24\SimpleUsableWidth}|K{0.76\SimpleUsableWidth}|}
  {\TableHeaderCell{Field} & \TableHeaderCell{Value}}
  Profile & \code{9e-rsa3072} \\\hline
  Source & Committed YubiKey loader material in test-vectors/piv-auto-yubikey-material/9e-rsa3072. Generated-card coverage. \\\hline
  Algorithm ID & \code{0x}\ApduHeader{05} \\\hline
  Key Reference & \code{0x}\ApduHeader{9E} \\\hline
  KMAC Challenge Digest & \ApduChallenge{4C98712B25DB7F477096FA3D0AC93FA2CA6BBE89CB291AC26EDF0D8312CA61D0CA238F1656247B2DF38E8B72D9AF7F5E} \\\hline
  Tag 81 Challenge Input & \ApduChallenge{0001FFFFFFFFFFFFFFFFFFFFFFFFFFFFFFFFFFFFFFFFFFFFFFFFFFFFFFFFFFFFFFFFFFFFFFFFFFFFFFFFFFFFFFFFFFFFFFFFFFFFFFFFFFFFFFFFFFFFFFFFFFFFFFFFFFFFFFFFFFFFFFFFFFFFFFFFFFFFFFFFFFFFFFFFFFFFFFFFFFFFFFFFFFFFFFFFFFFFFFFFFFFFFFFFFFFFFFFFFFFFFFFFFFFFFFFFFFFFFFFFFFFFFFFFFFFFFFFFFFFFFFFFFFFFFFFFFFFFFFFFFFFFFFFFFFFFFFFFFFFFFFFFFFFFFFFFFFFFFFFFFFFFFFFFFFFFFFFFFFFFFFFFFFFFFFFFFFFFFFFFFFFFFFFFFFFFFFFFFFFFFFFFFFFFFFFFFFFFFFFFFFFFFFFFFFFFFFFFFFFFFFFFFFFFFFFFFFFFFFFFFFFFFFFFFFFFFFFFFFFFFFFFFFFFFFFFFFFFFFFFFFFFFFFFFFFFFFFFFFFFFFFFFFFFFFFFFFFFFFFFFFFFFFFFFFFFFFFFFFFFFFFFFFFFFFFFFFFFFFFFFFFFFFFFFFFFFFFFFFFFFFFFFFFFFFFFFFFFFFFFFFFFFFFFFFFFFFFFFFFFFFFFFFFF003041300D0609608648016503040202050004304C98712B25DB7F477096FA3D0AC93FA2CA6BBE89CB291AC26EDF0D8312CA61D0CA238F1656247B2DF38E8B72D9AF7F5E} \\\hline
  GENERAL AUTHENTICATE APDU & \ApduHeader{0087059E}\ApduLength{00018A}\ApduTemplate{7C820186}\ApduResponse{8200}\ApduTemplate{81820180}\ApduChallenge{0001FFFFFFFFFFFFFFFFFFFFFFFFFFFFFFFFFFFFFFFFFFFFFFFFFFFFFFFFFFFFFFFFFFFFFFFFFFFFFFFFFFFFFFFFFFFFFFFFFFFFFFFFFFFFFFFFFFFFFFFFFFFFFFFFFFFFFFFFFFFFFFFFFFFFFFFFFFFFFFFFFFFFFFFFFFFFFFFFFFFFFFFFFFFFFFFFFFFFFFFFFFFFFFFFFFFFFFFFFFFFFFFFFFFFFFFFFFFFFFFFFFFFFFFFFFFFFFFFFFFFFFFFFFFFFFFFFFFFFFFFFFFFFFFFFFFFFFFFFFFFFFFFFFFFFFFFFFFFFFFFFFFFFFFFFFFFFFFFFFFFFFFFFFFFFFFFFFFFFFFFFFFFFFFFFFFFFFFFFFFFFFFFFFFFFFFFFFFFFFFFFFFFFFFFFFFFFFFFFFFFFFFFFFFFFFFFFFFFFFFFFFFFFFFFFFFFFFFFFFFFFFFFFFFFFFFFFFFFFFFFFFFFFFFFFFFFFFFFFFFFFFFFFFFFFFFFFFFFFFFFFFFFFFFFFFFFFFFFFFFFFFFFFFFFFFFFFFFFFFFFFFFFFFFFFFFFFFFFFFFFFFFFFFFFFFFFFFFFFFFFFFFFFFFFFFFFFFFFFFFFFFFFFFFF003041300D0609608648016503040202050004304C98712B25DB7F477096FA3D0AC93FA2CA6BBE89CB291AC26EDF0D8312CA61D0CA238F1656247B2DF38E8B72D9AF7F5E}\ApduLength{00} \\\hline
  Dynamic Authentication Response & \ApduTemplate{7C820184}\ApduResponse{82820180}\ApduResponse{74FA8827F5C2CFDAA5718679727015CD5A928828997F2F1423D033DE8AADEB3E86431170A714F5155B830F65F3F854120B16BC50F701C0EEA0278DB7EC2D7D6D587BEDB32F5E34EF1421E2C42543FF7DEEF7C5A41A01275B0BC126AF1208BBBD1E2C65149BC74C9E2CD00D822081B22E5CCB5090D2BD3970CA309D2124D79E9A28978D5DDF26E3FA8E37C360521C5BBE1AB27FACC24FF4C59AB46B7D571777B22E3E4981593D9F689DA3D33F3DE7C385C4315A7F154BE461968EA05A3315C4424717DD6F3CBF7FB94A135AE33573E58A18185F7E25771E6C4AB2B89D50007C7A60E6D65CF3BE58F8901B25B867D0FF97981BCBEE2BFCF5E57CEB5A81F40EEB44E52A0BD7F538F932AB4F1732828F2E95C4F203967ADF5983A285A127E6C0D7C2048261290B6B58D8110056DFE65547A7C1BDF78C6FEA62A8E07EC932AC6B58AD13C9423AAD5815C5251BCECEBBDB181F0B1C3621633734635FD48CD1AD16798C9B67CDF2BAB290DE0EA0BC5BA54E41736BA4885AF72D729A5168DC09DDD4B662}\ApduStatus{9000} \\\hline
  Response Tag 0x82 Length & 0x0180 bytes \\\hline
\end{SimpleTable}

\begin{SimpleTable}{PIV Auto Worked Example: PIV Authentication ECC P-256}{2}
  {|L{0.24\SimpleUsableWidth}|K{0.76\SimpleUsableWidth}|}
  {\TableHeaderCell{Field} & \TableHeaderCell{Value}}
  Profile & \code{9a-ecp256} \\\hline
  Source & Committed P-256 signing key in test-vectors/piv-auto-demo/source-p12, originally imported from GSA FICAM card-builder. NIST SD 33 / IR 8347 and generated-card coverage. \\\hline
  Algorithm ID & \code{0x}\ApduHeader{11} \\\hline
  Key Reference & \code{0x}\ApduHeader{9A} \\\hline
  KMAC Challenge Digest & \ApduChallenge{A3C81D3205FFB28A2AAAB0EDC21E8B3BEF5BD5F685FF5C686502CA5609B45851} \\\hline
  Tag 81 Challenge Input & \ApduChallenge{A3C81D3205FFB28A2AAAB0EDC21E8B3BEF5BD5F685FF5C686502CA5609B45851} \\\hline
  GENERAL AUTHENTICATE APDU & \ApduHeader{0087119A}\ApduLength{26}\ApduTemplate{7C24}\ApduResponse{8200}\ApduTemplate{8120}\ApduChallenge{A3C81D3205FFB28A2AAAB0EDC21E8B3BEF5BD5F685FF5C686502CA5609B45851}\ApduLength{00} \\\hline
  Dynamic Authentication Response & \ApduTemplate{7C4A}\ApduResponse{8248}\ApduResponse{30460221009D25A44E56572983819019DF26A2720E462A11D2CB3CD86A9142DEF4BCC1271B022100D084982BB9D3E2D3E6B0965956C71DBA0FE44FE328C0A2BD72951DFB4D9296FE}\ApduStatus{9000} \\\hline
  Response Tag 0x82 Length & 0x0048 bytes \\\hline
\end{SimpleTable}

\begin{SimpleTable}{PIV Auto Worked Example: Card Authentication ECC P-256}{2}
  {|L{0.24\SimpleUsableWidth}|K{0.76\SimpleUsableWidth}|}
  {\TableHeaderCell{Field} & \TableHeaderCell{Value}}
  Profile & \code{9e-ecp256} \\\hline
  Source & Committed P-256 signing key in test-vectors/piv-auto-demo/source-p12, originally imported from GSA FICAM card-builder. NIST SD 33 / IR 8347 and generated-card coverage. \\\hline
  Algorithm ID & \code{0x}\ApduHeader{11} \\\hline
  Key Reference & \code{0x}\ApduHeader{9E} \\\hline
  KMAC Challenge Digest & \ApduChallenge{965AF088B02DF65E7C299ABF53A5AB2DCBB3C789458DE8D0C819B19A9266658D} \\\hline
  Tag 81 Challenge Input & \ApduChallenge{965AF088B02DF65E7C299ABF53A5AB2DCBB3C789458DE8D0C819B19A9266658D} \\\hline
  GENERAL AUTHENTICATE APDU & \ApduHeader{0087119E}\ApduLength{26}\ApduTemplate{7C24}\ApduResponse{8200}\ApduTemplate{8120}\ApduChallenge{965AF088B02DF65E7C299ABF53A5AB2DCBB3C789458DE8D0C819B19A9266658D}\ApduLength{00} \\\hline
  Dynamic Authentication Response & \ApduTemplate{7C4A}\ApduResponse{8248}\ApduResponse{30460221008B204FA48FAA27A57D6E90F985199F284D2A8E0D0F2E051BBFBCACBA824619B9022100C6E04880D322DE7D10823EC78C8EA4CCB69FE394F694F1D12C797FA7BCAD7695}\ApduStatus{9000} \\\hline
  Response Tag 0x82 Length & 0x0048 bytes \\\hline
\end{SimpleTable}

\begin{SimpleTable}{PIV Auto Worked Example: PIV Authentication ECC P-384}{2}
  {|L{0.24\SimpleUsableWidth}|K{0.76\SimpleUsableWidth}|}
  {\TableHeaderCell{Field} & \TableHeaderCell{Value}}
  Profile & \code{9a-ecp384} \\\hline
  Source & Committed YubiKey loader material in test-vectors/piv-auto-yubikey-material/9a-ecp384. Generated-card coverage. \\\hline
  Algorithm ID & \code{0x}\ApduHeader{14} \\\hline
  Key Reference & \code{0x}\ApduHeader{9A} \\\hline
  KMAC Challenge Digest & \ApduChallenge{C813683447986D513FC17CED5DB4D335BFBEAFB1AF9E5200C6B7B1C1E3535D700C2E05F9011717E153998832B7588D28} \\\hline
  Tag 81 Challenge Input & \ApduChallenge{C813683447986D513FC17CED5DB4D335BFBEAFB1AF9E5200C6B7B1C1E3535D700C2E05F9011717E153998832B7588D28} \\\hline
  GENERAL AUTHENTICATE APDU & \ApduHeader{0087149A}\ApduLength{36}\ApduTemplate{7C34}\ApduResponse{8200}\ApduTemplate{8130}\ApduChallenge{C813683447986D513FC17CED5DB4D335BFBEAFB1AF9E5200C6B7B1C1E3535D700C2E05F9011717E153998832B7588D28}\ApduLength{00} \\\hline
  Dynamic Authentication Response & \ApduTemplate{7C6A}\ApduResponse{8268}\ApduResponse{3066023100D88ABABB9893D88341CFA2F1AB6E3D35E4AF02F4F8D085B5A3DE63574DCF694741D869066B8E6297429120614DCC9CE40231009A3F203BB5EA98839008B0D25F98C98DE8E87622F61279B2F2B29343912999A36311102B310F9985844BCC248350377B}\ApduStatus{9000} \\\hline
  Response Tag 0x82 Length & 0x0068 bytes \\\hline
\end{SimpleTable}

\begin{SimpleTable}{PIV Auto Worked Example: Card Authentication ECC P-384}{2}
  {|L{0.24\SimpleUsableWidth}|K{0.76\SimpleUsableWidth}|}
  {\TableHeaderCell{Field} & \TableHeaderCell{Value}}
  Profile & \code{9e-ecp384} \\\hline
  Source & Committed YubiKey loader material in test-vectors/piv-auto-yubikey-material/9e-ecp384. Generated-card coverage. \\\hline
  Algorithm ID & \code{0x}\ApduHeader{14} \\\hline
  Key Reference & \code{0x}\ApduHeader{9E} \\\hline
  KMAC Challenge Digest & \ApduChallenge{F3480C6C1DAD008D3E14D0C815D381F01420E9D3402A175BED097C6949E4BA447FA0A11745CE4D054A95B10A9D5CC689} \\\hline
  Tag 81 Challenge Input & \ApduChallenge{F3480C6C1DAD008D3E14D0C815D381F01420E9D3402A175BED097C6949E4BA447FA0A11745CE4D054A95B10A9D5CC689} \\\hline
  GENERAL AUTHENTICATE APDU & \ApduHeader{0087149E}\ApduLength{36}\ApduTemplate{7C34}\ApduResponse{8200}\ApduTemplate{8130}\ApduChallenge{F3480C6C1DAD008D3E14D0C815D381F01420E9D3402A175BED097C6949E4BA447FA0A11745CE4D054A95B10A9D5CC689}\ApduLength{00} \\\hline
  Dynamic Authentication Response & \ApduTemplate{7C69}\ApduResponse{8267}\ApduResponse{3065023034C6B2659FBE26A78D2615F8C0EBF2E9674292D9DE3CE8A31D9DA297B5AE236AC886E9E9C288D42474564BF0C46AE40C023100CCBD59A047D73F9FD58F60218FBCB4F99B18EBD7EB70EFE5CF64D5EB1A65F1676A867E1B1E5EB4CF9F34B44BE905D29C}\ApduStatus{9000} \\\hline
  Response Tag 0x82 Length & 0x0067 bytes \\\hline
\end{SimpleTable}


\subsection{APDUs}
\begin{SimpleTable}{PIV Auto APDU Reference: Card Preparation}{4}
  {|L{0.20\SimpleUsableWidth}|K{0.34\SimpleUsableWidth}|K{0.18\SimpleUsableWidth}|L{0.28\SimpleUsableWidth}|}
  {\TableHeaderCell{Operation} & \TableHeaderCell{Command APDU} & \TableHeaderCell{Success Response} & \TableHeaderCell{Purpose}}
  Select PIV Application & \ApduHeader{00A40400}\ApduLength{0B}\ApduTemplate{A000000308000010000100} & \ApduStatus{9000} & Selects the PIV application before reading objects or authenticating. \\\hline
  Get Discovery Object & \ApduHeader{00CB3FFF}\ApduLength{03}\ApduTemplate{5C017E} & \ApduTemplate{7E...}\ApduStatus{9000} & Reads the Discovery Object so the PD can determine application and PIN/VCI policy. \\\hline
  Verify PIN (optional) & \ApduHeader{00200080}\ApduLength{08}\ApduChallenge{313233343536FFFF} & \ApduStatus{9000} & Verifies example PIN \code{123456} padded with \ApduChallenge{FF}; failure responses include \ApduStatus{63CX} and \ApduStatus{6983}. \\\hline
\end{SimpleTable}

\begin{SimpleTable}{PIV Auto APDU Reference: GENERAL AUTHENTICATE Command}{4}
  {|L{0.18\SimpleUsableWidth}|K{0.28\SimpleUsableWidth}|L{0.18\SimpleUsableWidth}|L{0.36\SimpleUsableWidth}|}
  {\TableHeaderCell{Field} & \TableHeaderCell{Encoding} & \TableHeaderCell{Source} & \TableHeaderCell{Notes}}
  CLA & \ApduHeader{00} & ISO/IEC 7816-4 & Logical command before credential SM wrapping. Apply SP~800-73-5 Part~2, Section~4.2, when SM is established. \\\hline
  INS & \ApduHeader{87} & ISO/IEC 7816-4 & GENERAL AUTHENTICATE instruction. \\\hline
  P1 & Algorithm ID & Certificate SPKI & \ApduHeader{05}, \ApduHeader{07}, \ApduHeader{11}, or \ApduHeader{14} for supported PIV Auto profiles; \ApduHeader{06} appears only in informative RSA-1024 legacy examples. \\\hline
  P2 & Key Reference & Configuration Entry & \ApduHeader{9A} or \ApduHeader{9E} for PIV; \ApduHeader{9E} for TWIC Auto. \\\hline
  Lc & Short or extended & Command length & Short APDUs use one length byte; extended APDUs use \ApduLength{00} followed by two length bytes. \\\hline
  Command data & \ApduTemplate{7C...}\ApduResponse{8200}\ApduTemplate{81...}\ApduChallenge{...} & Dynamic Authentication Template & Tag \ApduResponse{82} is empty to request a response; tag \ApduTemplate{81} carries the exact challenge input. \\\hline
  Le & \ApduLength{00} & Response request & Requests the maximum response supported by the card interface. \\\hline
\end{SimpleTable}

\begin{SimpleTable}{PIV Auto APDU Reference: GENERAL AUTHENTICATE Response}{4}
  {|L{0.20\SimpleUsableWidth}|K{0.30\SimpleUsableWidth}|L{0.22\SimpleUsableWidth}|L{0.28\SimpleUsableWidth}|}
  {\TableHeaderCell{Case} & \TableHeaderCell{Response APDU} & \TableHeaderCell{PD Handling} & \TableHeaderCell{OSDP Payload}}
  Complete response & \ApduTemplate{7C...}\ApduResponse{82...}\ApduStatus{9000} & Strip only the ISO status word. & Carry the complete \ApduTemplate{7C} Dynamic Authentication Template in \texttt{osdp\_CRAUTHR} or \texttt{osdp\_PIVSTATUS}. \\\hline
  Additional response bytes & \ApduStatus{61XX} & Issue GET RESPONSE with Le \ApduLength{XX}. & Concatenate returned data before reporting the final \ApduTemplate{7C} template. \\\hline
  GET RESPONSE command & \ApduHeader{00C00000}\ApduLength{XX} & Retrieve remaining card response bytes. & Not exposed directly to the ACU. \\\hline
  Credential error & \ApduStatus{69..} or \ApduStatus{6A..} & Preserve command failure detail. & Auto reports terminal PIVSTATUS Result \code{0x06} with the applicable status word; commanded CRAUTH uses PIVERROR. \\\hline
\end{SimpleTable}

\begin{SimpleTable}{PIV Auto APDU Reference: OPACITY / VCI Context}{4}
  {|L{0.20\SimpleUsableWidth}|K{0.31\SimpleUsableWidth}|K{0.14\SimpleUsableWidth}|L{0.35\SimpleUsableWidth}|}
  {\TableHeaderCell{Operation} & \TableHeaderCell{Command Shape} & \TableHeaderCell{Key Ref} & \TableHeaderCell{Relationship to PIV Auto}}
  OPACITY establish & \ApduHeader{0087[27/2E]04}\ApduLength{[Lc]}\ApduTemplate{7C...81...}\ApduResponse{8200} & \ApduHeader{04} & Establishes credential SM after CVC and cryptogram verification. Contactless VCI also depends on Discovery policy and pairing when required. \\\hline
  PIV Auto challenge & \ApduHeader{0087[alg][9A/9E]}\ApduLength{[Lc]}\ApduTemplate{7C...}\ApduResponse{8200}\ApduTemplate{81...}\ApduChallenge{...}\ApduLength{00} & \ApduHeader{9A}/\ApduHeader{9E} & Performs the autonomous key challenge after card preparation and any required VCI setup. \\\hline
  Response handling & \ApduTemplate{7C...}\ApduResponse{82...}\ApduStatus{9000} & \ApduHeader{9A}/\ApduHeader{9E} & The PD reports the returned signed response bytes exactly, excluding only SW1/SW2. \\\hline
\end{SimpleTable}

\clearpage
\section{Credential SM and VCI Worked Examples (Informative)}
\label{app:sm-vci-examples}

These examples follow the preparation and PIN-verification flows in
Subsection~\ref{sec:credential-sm-flow}. A simulated ACU sends complete OSDP
packets to a simulated PD. The PD replays existing card captures and reports
the resulting credential state. These examples omit OSDP Secure Channel for
clarity; the test harness assumes that the OSDP Secure Channel authorization
requirements are satisfied.

The contact and contactless card01 captures supply the card commands and
replies. The verifier checks the CVC signature against an existing trust anchor,
recomputes SM key establishment, and checks protected command and response
MACs and encryption. The replay retains intervening captured exchanges so the
session counters and MAC chains remain intact. Those intervening exchanges
preserve the original card session. The card captures include earlier
provisioning and data-read commands; those commands impose no additional
requirements on automatic card preparation. The examples replay recorded
sessions; no connected card is used.

\begin{SimpleTable}{Replay coverage}{3}
  {|L{0.27\SimpleUsableWidth}|L{0.35\SimpleUsableWidth}|L{0.38\SimpleUsableWidth}|}
  {\TableHeaderCell{Case} & \TableHeaderCell{Evidence} & \TableHeaderCell{Purpose}}
  Contact SM & Captured SM session and CVC validation & Report SM with VCI absent. \\\hline
  Incorrect PIN & Modeled authenticated rejection & Clear the failed PIN while preserving SM and VCI. \\\hline
  Commanded pairing and PIN & Captured protected VERIFY exchanges & Complete pairing before verifying the cached PIN. \\\hline
  PIV Auto failure & Modeled results and full status words & Report the cause across multipart replies. \\\hline
  PIN cache expiration & Modeled timing, with captured VERIFY for the submitted-command case & Separate cache availability from credential PIN verification. \\\hline
\end{SimpleTable}

Complete packets, APDUs, and source references are in
\href{run:worked-examples/sm-vci/examples.json}{the SM/VCI example file}.
The accompanying README gives the replay and independent-verification commands.
Packet dumps use hexadecimal byte pairs, including the OSDP header and CRC.
Each example starts at sequence \code{0x01} on an established link to PD
address \code{0x01}. Standard status-query and pairing replies use the same
\texttt{osdp\_CARDSTATUSR} layout with Status Type \code{0x01}.
Diagnostic replies carry Status Type \code{0x02} before the multi-part fields
in every fragment. Timing examples identify the modeled elapsed time at each
exchange; captured APDUs retain their original order.

% Generated by scripts/render_sm_vci_examples.py; do not hand edit.
\Needspace{24\baselineskip}
\subsection{Credential SM over contact}
\label{ex:sm-contact-sm}
The PD prepares a contact credential and reports SM established without VCI. The reply identifies the CVC issuer and trust anchor used for that SM session.

Example ID: \texttt{contact-sm}.
Credential exchanges come from the contact card01 capture. The complete replay retains its original APDU order.
\begin{SimpleTable}{Reported security state}{3}
{|L{0.30\SimpleUsableWidth}|L{0.32\SimpleUsableWidth}|L{0.38\SimpleUsableWidth}|}
{\TableHeaderCell{Field} & \TableHeaderCell{Value} & \TableHeaderCell{Notes}}
Status Type & \code{0x01} & Standard credential status. \\\hline
Credential Security Status & \code{0x02} & SM established without VCI. \\\hline
SM Trust Anchor ID & \code{0x01} & Loaded anchor used to validate SM. \\\hline
SM CVC IIN & \code{0xB6103792270FBD08} & Issuer named by the credential's SM CVC. \\\hline
SM Anchor IIN & \code{0xB6103792270FBD08} & Same issuer: direct trust path. \\\hline
\end{SimpleTable}
\begingroup\small
\begin{SimpleTable}{OSDP command and reply exchange}{3}
{|L{0.20\SimpleUsableWidth}|L{0.50\SimpleUsableWidth}|L{0.30\SimpleUsableWidth}|}
{\TableHeaderCell{Message} & \TableHeaderCell{Packet bytes (hex)} & \TableHeaderCell{Notes}}
ACU\newline \texttt{CARDSTATUS} & \texttt{53 01 09 00 05 A9 01 8C 9B} & Read current state after automatic credential preparation. \\\hline
PD\newline \texttt{ACK} & \texttt{53 81 08 00 05 40 39 35} & Command accepted. \\\hline
ACU\newline \texttt{POLL} & \texttt{53 01 08 00 06 60 D8 66} & Poll for completion. \\\hline
PD\newline \texttt{CARDSTATUSR} & \texttt{53 81 2A 00 06 85 01 01 02 01 B6 10}\newline \texttt{37 92 27 0F BD 08 B6 10 37 92 27 0F}\newline \texttt{BD 08 48 00 0B A0 00 00 03 08 00 00}\newline \texttt{10 00 01 00 98 79} & Status 0x02: SM established without VCI. \\\hline
\end{SimpleTable}
\endgroup

\Needspace{24\baselineskip}
\subsection{PIN before credential presentation}
\label{ex:sm-cache-before-card}
The PD collects the PIN locally and reports completion before the credential is presented. Initial selection and SM establishment preserve the cache. The ACU then completes pairing and requests PIN verification over VCI.

Example ID: \texttt{cache-before-card}.
Credential exchanges come from the contactless card01 capture. The complete replay retains its original APDU order.
\begingroup\small
\begin{SimpleTable}{OSDP command and reply exchange}{3}
{|L{0.20\SimpleUsableWidth}|L{0.50\SimpleUsableWidth}|L{0.30\SimpleUsableWidth}|}
{\TableHeaderCell{Message} & \TableHeaderCell{Packet bytes (hex)} & \TableHeaderCell{Notes}}
ACU\newline \texttt{PIVSPE} & \texttt{53 01 0E 00 05 AA 04 00 00 00 3C 00}\newline \texttt{FE 9F} & Collect a PIN locally; cache for 0x003C seconds. \\\hline
PD\newline \texttt{ACK} & \texttt{53 81 08 00 05 40 39 35} & Command accepted. \\\hline
ACU\newline \texttt{POLL} & \texttt{53 01 08 00 06 60 D8 66} & Poll for completion. \\\hline
PD\newline \texttt{PIVSPER} & \texttt{53 81 09 00 06 86 00 05 10} & Result 0x00: requested PIN operation completed. \\\hline
ACU\newline \texttt{CARDSTATUS} & \texttt{53 01 09 00 07 A9 01 EC F5} & Read current state after automatic credential preparation. \\\hline
PD\newline \texttt{ACK} & \texttt{53 81 08 00 07 40 5B 53} & Command accepted. \\\hline
ACU\newline \texttt{POLL} & \texttt{53 01 08 00 05 60 8B 33} & Poll for completion. \\\hline
PD\newline \texttt{CARDSTATUSR} & \texttt{53 81 2A 00 05 85 01 02 02 01 B6 10}\newline \texttt{37 92 27 0F BD 08 B6 10 37 92 27 0F}\newline \texttt{BD 08 48 00 0B A0 00 00 03 08 00 00}\newline \texttt{10 00 01 00 1E 84} & Status 0x02: SM established without VCI. \\\hline
ACU\newline \texttt{PIVXMITPAIRING} & \texttt{53 01 10 00 06 AD 30 30 30 30 30 30}\newline \texttt{30 32 EB 8B} & Submit the captured pairing code. \\\hline
PD\newline \texttt{ACK} & \texttt{53 81 08 00 06 40 6A 60} & Command accepted. \\\hline
ACU\newline \texttt{POLL} & \texttt{53 01 08 00 07 60 E9 55} & Poll for completion. \\\hline
PD\newline \texttt{CARDSTATUSR} & \texttt{53 81 2A 00 07 85 01 02 01 01 B6 10}\newline \texttt{37 92 27 0F BD 08 B6 10 37 92 27 0F}\newline \texttt{BD 08 48 00 0B A0 00 00 03 08 00 00}\newline \texttt{10 00 01 00 56 3D} & Status 0x01: SM and VCI established. \\\hline
ACU\newline \texttt{PIVSPE} & \texttt{53 01 0E 00 05 AA 02 00 00 00 00 00}\newline \texttt{E7 52} & Verify the cached PIN over established VCI. \\\hline
PD\newline \texttt{ACK} & \texttt{53 81 08 00 05 40 39 35} & Command accepted. \\\hline
ACU\newline \texttt{POLL} & \texttt{53 01 08 00 06 60 D8 66} & Poll for completion. \\\hline
PD\newline \texttt{PIVSPER} & \texttt{53 81 09 00 06 86 00 05 10} & Result 0x00: requested PIN operation completed. \\\hline
\end{SimpleTable}
\endgroup

\Needspace{24\baselineskip}
\subsection{PIV Auto needs commanded pairing}
\label{ex:sm-auto-pairing-required}
In this modeled case, the configured PIV Auto authentication requires VCI. SM is established, but pairing remains necessary. PIV Auto reports the unmet condition and ends the attempt. A later poll receives ACK; the ACU can begin commanded pairing.

Example ID: \texttt{auto-pairing-required}.
\begingroup\small
\begin{SimpleTable}{OSDP command and reply exchange}{3}
{|L{0.20\SimpleUsableWidth}|L{0.50\SimpleUsableWidth}|L{0.30\SimpleUsableWidth}|}
{\TableHeaderCell{Message} & \TableHeaderCell{Packet bytes (hex)} & \TableHeaderCell{Notes}}
ACU\newline \texttt{POLL} & \texttt{53 01 08 00 05 60 8B 33} & Poll for completion. \\\hline
PD\newline \texttt{PIVSTATUS} & \texttt{53 81 45 00 05 89 37 00 00 00 37 00}\newline \texttt{05 01 01 00 00 00 02 00 00 00 00 00}\newline \texttt{00 00 00 00 00 00 00 00 00 00 00 00}\newline \texttt{00 00 00 00 00 00 00 00 00 00 00 00}\newline \texttt{00 00 00 00 00 00 00 00 00 00 00 00}\newline \texttt{00 00 00 9A 07 00 00 FF 09} & Result 0x05, Detail 0x01: pairing code required. \\\hline
ACU\newline \texttt{POLL} & \texttt{53 01 08 00 06 60 D8 66} & Poll for completion. \\\hline
PD\newline \texttt{ACK} & \texttt{53 81 08 00 06 40 6A 60} & No completion pending. \\\hline
\end{SimpleTable}
\endgroup

\Needspace{24\baselineskip}
\subsection{SM failure diagnostics}
\label{ex:sm-diagnostic-cvc-hash}
In this modeled case, PIV Auto reports Result 0x02, Detail 0x04: invalid authentication cryptogram. The ACU then requests the SM diagnostic record. The PD returns the failure reason and the hash of the received CVC, with no further credential commands. Three fragments carry the 35-byte record. The CVC bytes used for the hash come from the existing capture.

Example ID: \texttt{diagnostic-cvc-hash}.
\begingroup\small
\begin{SimpleTable}{OSDP command and reply exchange}{3}
{|L{0.20\SimpleUsableWidth}|L{0.50\SimpleUsableWidth}|L{0.30\SimpleUsableWidth}|}
{\TableHeaderCell{Message} & \TableHeaderCell{Packet bytes (hex)} & \TableHeaderCell{Notes}}
ACU\newline \texttt{CARDSTATUS} & \texttt{53 01 09 00 06 A9 02 BF F2} & Request retained SM diagnostics; no credential command is sent. \\\hline
PD\newline \texttt{ACK} & \texttt{53 81 08 00 06 40 6A 60} & Command accepted. \\\hline
ACU\newline \texttt{POLL} & \texttt{53 01 08 00 07 60 E9 55} & Poll for completion. \\\hline
PD\newline \texttt{CARDSTATUSR} & \texttt{53 81 1F 00 07 85 02 23 00 00 00 10}\newline \texttt{00 03 04 01 A5 50 00 73 BA 17 70 37}\newline \texttt{0A 1D 58 B1 5B 2F F2} & Diagnostic: total 0x0023, offset 0x0000, fragment 0x0010. \\\hline
ACU\newline \texttt{POLL} & \texttt{53 01 08 00 05 60 8B 33} & Poll for completion. \\\hline
PD\newline \texttt{CARDSTATUSR} & \texttt{53 81 1F 00 05 85 02 23 00 10 00 10}\newline \texttt{00 97 5D D1 50 7F 39 22 BB 44 20 69}\newline \texttt{9B 2D FB 42 83 52 B1} & Diagnostic: total 0x0023, offset 0x0010, fragment 0x0010. \\\hline
ACU\newline \texttt{POLL} & \texttt{53 01 08 00 06 60 D8 66} & Poll for completion. \\\hline
PD\newline \texttt{CARDSTATUSR} & \texttt{53 81 12 00 06 85 02 23 00 20 00 03}\newline \texttt{00 C7 6C E9 B4 96} & Diagnostic: total 0x0023, offset 0x0020, fragment 0x0003. \\\hline
\end{SimpleTable}
\endgroup

\Needspace{24\baselineskip}
\subsection{PIN cache expires while idle}
\label{ex:sm-pin-expiration-idle}
The PD acquires a PIN at t=0 with a ten-second validity period. At t=10 the PD erases the PIN. A poll receives ACK because no operation is pending. A status query reports an empty cache; a subsequent verification request ends with no cached PIN. The PD sends no VERIFY. This example models the timing and credential state.

Example ID: \texttt{pin-expiration-idle}.
\label{ex:pin-expiration}
\begingroup\small
\begin{SimpleTable}{OSDP command and reply exchange}{3}
{|L{0.20\SimpleUsableWidth}|L{0.50\SimpleUsableWidth}|L{0.30\SimpleUsableWidth}|}
{\TableHeaderCell{Message} & \TableHeaderCell{Packet bytes (hex)} & \TableHeaderCell{Notes}}
ACU\newline \texttt{PIVSPE} & \texttt{53 01 0E 00 05 AA 04 00 00 00 0A 00}\newline \texttt{CD 30} & t=0 s. Acquire a PIN; Cache Validity is 0x000A seconds. \\\hline
PD\newline \texttt{ACK} & \texttt{53 81 08 00 05 40 39 35} & t=0 s. Command accepted. \\\hline
ACU\newline \texttt{POLL} & \texttt{53 01 08 00 06 60 D8 66} & t=0 s. Poll for completion. \\\hline
PD\newline \texttt{PIVSPER} & \texttt{53 81 09 00 06 86 00 05 10} & t=0 s. Result 0x00: requested PIN operation completed. \\\hline
ACU\newline \texttt{POLL} & \texttt{53 01 08 00 07 60 E9 55} & t=10 s. Poll for completion. \\\hline
PD\newline \texttt{ACK} & \texttt{53 81 08 00 07 40 5B 53} & t=10 s. No completion pending. \\\hline
ACU\newline \texttt{PIVSPE} & \texttt{53 01 0E 00 05 AA 00 00 00 00 00 00}\newline \texttt{A7 D9} & t=10 s. Report empty cache and no PIN verification. \\\hline
PD\newline \texttt{PIVSPER} & \texttt{53 81 09 00 05 86 00 55 49} & t=10 s. Cache empty; credential PIN verification absent. \\\hline
ACU\newline \texttt{PIVSPE} & \texttt{53 01 0E 00 06 AA 03 00 00 00 00 00}\newline \texttt{32 DF} & t=10 s. Request verification using the expired cache. \\\hline
PD\newline \texttt{ACK} & \texttt{53 81 08 00 06 40 6A 60} & t=10 s. Command accepted. \\\hline
ACU\newline \texttt{POLL} & \texttt{53 01 08 00 07 60 E9 55} & t=10 s. Poll for completion. \\\hline
PD\newline \texttt{PIVSPER} & \texttt{53 81 09 00 07 86 0A 7F 86} & t=10 s. Result 0x0A: no cached PIN. \\\hline
\end{SimpleTable}
\endgroup

\Needspace{24\baselineskip}
\subsection{Expiration while waiting for a credential}
\label{ex:sm-pin-expiration-waiting}
The ACU arms cached-PIN verification with Auto Execute flag 0x02. The credential is absent when the cache expires at t=10. The PD ends the pending request with PIVSPER 0x0A (no cached PIN). This example models the waiting state.

Example ID: \texttt{pin-expiration-waiting}.
\begingroup\small
\begin{SimpleTable}{OSDP command and reply exchange}{3}
{|L{0.20\SimpleUsableWidth}|L{0.50\SimpleUsableWidth}|L{0.30\SimpleUsableWidth}|}
{\TableHeaderCell{Message} & \TableHeaderCell{Packet bytes (hex)} & \TableHeaderCell{Notes}}
ACU\newline \texttt{PIVSPE} & \texttt{53 01 0E 00 05 AA 04 00 00 00 0A 00}\newline \texttt{CD 30} & t=0 s. Acquire a PIN; Cache Validity is 0x000A seconds. \\\hline
PD\newline \texttt{ACK} & \texttt{53 81 08 00 05 40 39 35} & t=0 s. Command accepted. \\\hline
ACU\newline \texttt{POLL} & \texttt{53 01 08 00 06 60 D8 66} & t=0 s. Poll for completion. \\\hline
PD\newline \texttt{PIVSPER} & \texttt{53 81 09 00 06 86 00 05 10} & t=0 s. Result 0x00: requested PIN operation completed. \\\hline
ACU\newline \texttt{PIVSPE} & \texttt{53 01 0E 00 07 AA 03 00 00 02 00 00}\newline \texttt{81 F6} & t=0 s. Arm cached-PIN verification for credential presentation. \\\hline
PD\newline \texttt{ACK} & \texttt{53 81 08 00 07 40 5B 53} & t=0 s. Command accepted. \\\hline
ACU\newline \texttt{POLL} & \texttt{53 01 08 00 05 60 8B 33} & t=0 s. Poll for completion. \\\hline
PD\newline \texttt{ACK} & \texttt{53 81 08 00 05 40 39 35} & t=0 s. No completion pending. \\\hline
ACU\newline \texttt{POLL} & \texttt{53 01 08 00 06 60 D8 66} & t=10 s. Poll for completion. \\\hline
PD\newline \texttt{PIVSPER} & \texttt{53 81 09 00 06 86 0A 4F B1} & t=10 s. Result 0x0A: no cached PIN. \\\hline
\end{SimpleTable}
\endgroup

\Needspace{24\baselineskip}
\subsection{Expiration before VERIFY submission}
\label{ex:sm-pin-expiration-boundary}
The PD accepts verification at t=9. Before the PD submits VERIFY, the deadline is reached at t=10. The PD returns PIVSPER 0x0A (no cached PIN) and sends no VERIFY. Acceptance alone does not reserve the PIN beyond its deadline. Timing is modeled.

Example ID: \texttt{pin-expiration-boundary}.
\begingroup\small
\begin{SimpleTable}{OSDP command and reply exchange}{3}
{|L{0.20\SimpleUsableWidth}|L{0.50\SimpleUsableWidth}|L{0.30\SimpleUsableWidth}|}
{\TableHeaderCell{Message} & \TableHeaderCell{Packet bytes (hex)} & \TableHeaderCell{Notes}}
ACU\newline \texttt{PIVSPE} & \texttt{53 01 0E 00 05 AA 04 00 00 00 0A 00}\newline \texttt{CD 30} & t=0 s. Acquire a PIN; Cache Validity is 0x000A seconds. \\\hline
PD\newline \texttt{ACK} & \texttt{53 81 08 00 05 40 39 35} & t=0 s. Command accepted. \\\hline
ACU\newline \texttt{POLL} & \texttt{53 01 08 00 06 60 D8 66} & t=0 s. Poll for completion. \\\hline
PD\newline \texttt{PIVSPER} & \texttt{53 81 09 00 06 86 00 05 10} & t=0 s. Result 0x00: requested PIN operation completed. \\\hline
ACU\newline \texttt{PIVSPE} & \texttt{53 01 0E 00 07 AA 03 00 00 00 00 00}\newline \texttt{E1 98} & t=9 s. Accept verification at t=9, before the cache deadline. \\\hline
PD\newline \texttt{ACK} & \texttt{53 81 08 00 07 40 5B 53} & t=9 s. Command accepted. \\\hline
ACU\newline \texttt{POLL} & \texttt{53 01 08 00 05 60 8B 33} & t=10 s. Poll for completion. \\\hline
PD\newline \texttt{PIVSPER} & \texttt{53 81 09 00 05 86 0A 1F E8} & t=10 s. Result 0x0A: no cached PIN. \\\hline
\end{SimpleTable}
\endgroup

\Needspace{24\baselineskip}
\subsection{Expiration after VERIFY submission}
\label{ex:sm-pin-expiration-submitted}
The PD submits VERIFY at t=9, before the deadline. The cache expires at t=10 while the exchange is completing. The PD reports successful verification, then reports an empty cache with the credential still PIN verified. Credential APDUs come from the contactless capture; the timing is modeled.

Example ID: \texttt{pin-expiration-submitted}.
\begingroup\small
\begin{SimpleTable}{OSDP command and reply exchange}{3}
{|L{0.20\SimpleUsableWidth}|L{0.50\SimpleUsableWidth}|L{0.30\SimpleUsableWidth}|}
{\TableHeaderCell{Message} & \TableHeaderCell{Packet bytes (hex)} & \TableHeaderCell{Notes}}
ACU\newline \texttt{PIVSPE} & \texttt{53 01 0E 00 05 AA 04 00 00 00 0A 00}\newline \texttt{CD 30} & t=0 s. Acquire a PIN; Cache Validity is 0x000A seconds. \\\hline
PD\newline \texttt{ACK} & \texttt{53 81 08 00 05 40 39 35} & t=0 s. Command accepted. \\\hline
ACU\newline \texttt{POLL} & \texttt{53 01 08 00 06 60 D8 66} & t=0 s. Poll for completion. \\\hline
PD\newline \texttt{PIVSPER} & \texttt{53 81 09 00 06 86 00 05 10} & t=0 s. Result 0x00: requested PIN operation completed. \\\hline
ACU\newline \texttt{PIVXMITPAIRING} & \texttt{53 01 10 00 07 AD 30 30 30 30 30 30}\newline \texttt{30 32 AE E4} & t=0 s. Supply the captured pairing code after SM preparation. \\\hline
PD\newline \texttt{ACK} & \texttt{53 81 08 00 07 40 5B 53} & t=0 s. Command accepted. \\\hline
ACU\newline \texttt{POLL} & \texttt{53 01 08 00 05 60 8B 33} & t=0 s. Poll for completion. \\\hline
PD\newline \texttt{CARDSTATUSR} & \texttt{53 81 2A 00 05 85 01 02 01 01 B6 10}\newline \texttt{37 92 27 0F BD 08 B6 10 37 92 27 0F}\newline \texttt{BD 08 48 00 0B A0 00 00 03 08 00 00}\newline \texttt{10 00 01 00 D4 6F} & t=0 s. Status 0x01: SM and VCI established. \\\hline
ACU\newline \texttt{PIVSPE} & \texttt{53 01 0E 00 06 AA 02 00 00 00 00 00}\newline \texttt{92 9A} & t=9 s. Accept verification at t=9, before the cache deadline. \\\hline
PD\newline \texttt{ACK} & \texttt{53 81 08 00 06 40 6A 60} & t=9 s. Command accepted. \\\hline
ACU\newline \texttt{POLL} & \texttt{53 01 08 00 07 60 E9 55} & t=10 s. Poll for completion. \\\hline
PD\newline \texttt{PIVSPER} & \texttt{53 81 09 00 07 86 00 35 27} & t=10 s. Result 0x00: requested PIN operation completed. \\\hline
ACU\newline \texttt{PIVSPE} & \texttt{53 01 0E 00 05 AA 00 00 00 00 00 00}\newline \texttt{A7 D9} & t=10 s. Empty cache; credential PIN verified. \\\hline
PD\newline \texttt{PIVSPER} & \texttt{53 81 09 00 05 86 02 17 69} & t=10 s. Cache empty; credential PIN verified. \\\hline
\end{SimpleTable}
\endgroup

\Needspace{24\baselineskip}
\subsection{PIN reuse preserves the original deadline}
\label{ex:sm-pin-expiration-reuse}
Two explicit contact verification requests use the same cached PIN at t=1 and t=2. Both succeed. At t=10, the cache expires and the credential remains PIN verified. The successes and timing are modeled; the example demonstrates the PD state rules.

Example ID: \texttt{pin-expiration-reuse}.
\begingroup\small
\begin{SimpleTable}{OSDP command and reply exchange}{3}
{|L{0.20\SimpleUsableWidth}|L{0.50\SimpleUsableWidth}|L{0.30\SimpleUsableWidth}|}
{\TableHeaderCell{Message} & \TableHeaderCell{Packet bytes (hex)} & \TableHeaderCell{Notes}}
ACU\newline \texttt{PIVSPE} & \texttt{53 01 0E 00 05 AA 04 00 00 00 0A 00}\newline \texttt{CD 30} & t=0 s. Acquire a PIN; Cache Validity is 0x000A seconds. \\\hline
PD\newline \texttt{ACK} & \texttt{53 81 08 00 05 40 39 35} & t=0 s. Command accepted. \\\hline
ACU\newline \texttt{POLL} & \texttt{53 01 08 00 06 60 D8 66} & t=0 s. Poll for completion. \\\hline
PD\newline \texttt{PIVSPER} & \texttt{53 81 09 00 06 86 00 05 10} & t=0 s. Result 0x00: requested PIN operation completed. \\\hline
ACU\newline \texttt{PIVSPE} & \texttt{53 01 0E 00 07 AA 03 00 00 00 00 00}\newline \texttt{E1 98} & t=1 s. Explicit contact verification using the same cached PIN. \\\hline
PD\newline \texttt{ACK} & \texttt{53 81 08 00 07 40 5B 53} & t=1 s. Command accepted. \\\hline
ACU\newline \texttt{POLL} & \texttt{53 01 08 00 05 60 8B 33} & t=1 s. Poll for completion. \\\hline
PD\newline \texttt{PIVSPER} & \texttt{53 81 09 00 05 86 00 55 49} & t=1 s. Result 0x00: requested PIN operation completed. \\\hline
ACU\newline \texttt{PIVSPE} & \texttt{53 01 0E 00 06 AA 03 00 00 00 00 00}\newline \texttt{32 DF} & t=2 s. Explicit contact verification using the same cached PIN. \\\hline
PD\newline \texttt{ACK} & \texttt{53 81 08 00 06 40 6A 60} & t=2 s. Command accepted. \\\hline
ACU\newline \texttt{POLL} & \texttt{53 01 08 00 07 60 E9 55} & t=2 s. Poll for completion. \\\hline
PD\newline \texttt{PIVSPER} & \texttt{53 81 09 00 07 86 00 35 27} & t=2 s. Result 0x00: requested PIN operation completed. \\\hline
ACU\newline \texttt{PIVSPE} & \texttt{53 01 0E 00 05 AA 00 00 00 00 00 00}\newline \texttt{A7 D9} & t=10 s. Cache expired; credential PIN verification remains valid. \\\hline
PD\newline \texttt{PIVSPER} & \texttt{53 81 09 00 05 86 02 17 69} & t=10 s. Cache empty; credential PIN verified. \\\hline
\end{SimpleTable}
\endgroup

\Needspace{24\baselineskip}
\subsection{Credential failure status across two fragments}
\label{ex:sm-auto-status-word}
PIV Auto reports a modeled GENERAL AUTHENTICATE rejection: Result 0x06, Detail 0x00. The credential status is 0x6982 (security status not satisfied). The first fragment ends with SW1, 0x69; the second carries SW2, 0x82. The ACU reassembles the complete record before interpreting the status word.

Example ID: \texttt{auto-status-word}.
\begin{figure}[htbp]
\centering
\begin{tikzpicture}[font=\sffamily,
  record/.style={draw=OSDPBlue,fill=OSDPLightGray,text width=12cm,align=center,inner sep=9pt},
  flow/.style={->,draw=OSDPBlue,thick}]
\node[record] (logical) {\textbf{Complete logical reply: 0x0039 bytes}\\[3pt]
0x0037 fixed bytes, including Signed Response Length = 0\\
followed by Status Word: \texttt{69 82}};
\node[record,anchor=north] (first) at ([yshift=-0.65cm]logical.south)
{\textbf{First POLL: PIVSTATUS fragment}\\[3pt]
Total 0x0039; offset 0x0000; size 0x0038\\
Fixed fields and SW1 (\texttt{69})};
\node[record,anchor=north] (second) at ([yshift=-0.5cm]first.south)
{\textbf{Next POLL: PIVSTATUS fragment}\\[3pt]
Total 0x0039; offset 0x0038; size 0x0001\\
SW2 (\texttt{82})};
\node[record,anchor=north] (decode) at ([yshift=-0.65cm]second.south)
{\textbf{ACU reassembles, then interprets}\\[3pt]
Result 0x06: credential authentication command failed\\
Status 0x6982: credential security status not satisfied};
\draw[flow] (logical.south) -- (first.north);
\draw[flow] (first.south) -- (second.north);
\draw[flow] (second.south) -- (decode.north);
\end{tikzpicture}
\caption{OSDP 3.0 multipart fields describe the entire logical reply.
A fragment boundary can fall between SW1 and SW2.}
\label{fig:auto-status-word}
\end{figure}

\begingroup\small
\begin{SimpleTable}{OSDP command and reply exchange}{3}
{|L{0.20\SimpleUsableWidth}|L{0.50\SimpleUsableWidth}|L{0.30\SimpleUsableWidth}|}
{\TableHeaderCell{Message} & \TableHeaderCell{Packet bytes (hex)} & \TableHeaderCell{Notes}}
ACU\newline \texttt{POLL} & \texttt{53 01 08 00 05 60 8B 33} & Poll for completion. \\\hline
PD\newline \texttt{PIVSTATUS} & \texttt{53 81 46 00 05 89 39 00 00 00 38 00}\newline \texttt{06 00 01 00 00 00 02 00 00 00 00 00}\newline \texttt{00 00 00 00 00 00 00 00 00 00 00 00}\newline \texttt{00 00 00 00 00 00 00 00 00 00 00 00}\newline \texttt{00 00 00 00 00 00 00 00 00 00 00 00}\newline \texttt{00 00 00 9A 07 00 00 69 E4 69} & PIV Auto credential rejection; SW1/SW2 cross the fragment boundary. \\\hline
ACU\newline \texttt{POLL} & \texttt{53 01 08 00 06 60 D8 66} & Poll for completion. \\\hline
PD\newline \texttt{PIVSTATUS} & \texttt{53 81 0F 00 06 89 39 00 38 00 01 00}\newline \texttt{82 9D EB} & PIV Auto credential rejection; SW1/SW2 cross the fragment boundary. \\\hline
\end{SimpleTable}
\endgroup

\Needspace{24\baselineskip}
\subsection{Incorrect PIN preserves SM and VCI}
\label{ex:sm-sm-credential-rejection}
This modeled case starts with SM, VCI, and a cached PIN available. The credential returns an authenticated 0x63C2 (incorrect PIN, two attempts remain). The PD reports PIN mismatch and clears the cache and PIN-verification state. SM and VCI remain established. The SPE reply carries the mismatch result; a separate Mode 0x06 query retrieves remaining attempts. The ACU can acquire a corrected PIN and request verification again using the existing session.

Example ID: \texttt{sm-credential-rejection}.
\begingroup\small
\begin{SimpleTable}{OSDP command and reply exchange}{3}
{|L{0.20\SimpleUsableWidth}|L{0.50\SimpleUsableWidth}|L{0.30\SimpleUsableWidth}|}
{\TableHeaderCell{Message} & \TableHeaderCell{Packet bytes (hex)} & \TableHeaderCell{Notes}}
ACU\newline \texttt{PIVSPE} & \texttt{53 01 0E 00 05 AA 02 00 00 00 00 00}\newline \texttt{E7 52} & Verify a cached PIN over an established VCI; initial state is modeled. \\\hline
PD\newline \texttt{ACK} & \texttt{53 81 08 00 05 40 39 35} & Command accepted. \\\hline
ACU\newline \texttt{POLL} & \texttt{53 01 08 00 06 60 D8 66} & Poll for completion. \\\hline
PD\newline \texttt{PIVSPER} & \texttt{53 81 09 00 06 86 03 66 20} & Result 0x03: PIN mismatch. \\\hline
ACU\newline \texttt{PIVSPE} & \texttt{53 01 0E 00 07 AA 00 00 00 00 00 00}\newline \texttt{01 56} & Cache erased and PIN-verification state cleared after PIN mismatch. \\\hline
PD\newline \texttt{PIVSPER} & \texttt{53 81 09 00 07 86 00 35 27} & Cache empty; credential PIN verification absent. \\\hline
ACU\newline \texttt{CARDSTATUS} & \texttt{53 01 09 00 05 A9 01 8C 9B} & Read credential status after ordinary rejection. \\\hline
PD\newline \texttt{ACK} & \texttt{53 81 08 00 05 40 39 35} & Command accepted. \\\hline
ACU\newline \texttt{POLL} & \texttt{53 01 08 00 06 60 D8 66} & Poll for completion. \\\hline
PD\newline \texttt{CARDSTATUSR} & \texttt{53 81 2A 00 06 85 01 02 01 01 00 00}\newline \texttt{00 00 00 00 00 00 00 00 00 00 00 00}\newline \texttt{00 00 FF FF 0B A0 00 00 03 08 00 00}\newline \texttt{10 00 01 00 A9 03} & Status 0x01: SM and VCI established. \\\hline
\end{SimpleTable}
\endgroup



\subsection{Recovering trust for contactless PIV Auto}
\label{ex:sm-recovery}

This example uses a credential that supports SM and requires a pairing code
for VCI. PIV Auto is configured for contactless authentication with key
\code{0x9A}, but the PD's trust store is empty. The ACU already holds the roots
and validation policy authorized for the deployment. Figure~\ref{fig:sm-recovery}
shows the recovery sequence.

\begin{figure}[htbp]
\centering
\begin{tikzpicture}[x=1cm,y=1cm,font=\sffamily\small,
  recoverybox/.style={draw=OSDPBlue,rounded corners=3pt,fill=OSDPLightGray,
    text width=12.5cm,align=left,inner sep=9pt},
  flow/.style={->,draw=OSDPBlue,thick}]
\node[recoverybox] (failure) at (0,0) {
  \textbf{1. PD reports why PIV Auto stopped}\\[3pt]
  This PIN-based contactless 0x9A workflow requires SM, VCI, and a verified PIN.
  An empty trust store produces PIVSTATUS Result 0x02, Detail 0x01.
  The PD ends this attempt.};
\node[recoverybox,anchor=north] (inspect) at ([yshift=-0.45cm]failure.south) {
  \textbf{2. ACU obtains the diagnostic and signer information}\\[3pt]
  CARDSTATUS Type 0x02 requests CARDSTATUSR Type 0x02 fragments containing
  the setup outcome and any received CVC's hash.
  PIVGETDATA for 0x5FC122 reads the signer object from the credential.
  Recovery stops if that read fails. An empty store means setup was skipped, so the
  diagnostic reports that no CVC hash is available.};
\node[recoverybox,anchor=north] (trust) at ([yshift=-0.45cm]inspect.south) {
  \textbf{3. ACU validates and loads the trust anchor}\\[3pt]
  The ACU validates the signer certificate against an authorized root and
  deployment policy. After validation, PIVVCILOADTA loads the anchor.
  The ACU reapplies PIVMODE to start a new automatic attempt.};
\node[recoverybox,anchor=north] (pair) at ([yshift=-0.45cm]trust.south) {
  \textbf{4. PD reports pairing required; ACU supplies the code}\\[3pt]
  SM succeeds, but this credential requires pairing. PIV Auto reports
  Result 0x05, Detail 0x01 and ends the attempt. The ACU sends PIVXMITPAIRING
  and waits for status confirming SM and VCI.};
\node[recoverybox,anchor=north] (auth) at ([yshift=-0.45cm]pair.south) {
  \textbf{5. ACU completes authentication in commanded mode}\\[3pt]
  PIVSPE Mode 0x01 collects and verifies the PIN. The ACU obtains the 0x9A
  certificate with PIVGETDATA and sends a fresh challenge with CRAUTH.
  The ACU validates the certificate and response before deciding access.
  The existing SM and PIN state remain in use.};
\draw[flow] (failure.south) -- (inspect.north);
\draw[flow] (inspect.south) -- (trust.north);
\draw[flow] (trust.south) -- (pair.north);
\draw[flow] (pair.south) -- (auth.north);
\end{tikzpicture}
\caption{Recovery from an empty PD trust store for contactless PIV authentication
with key 0x9A. Every step depends on a successful result from the preceding step.
Unsupported SM or VCI requires a different ACU-selected workflow.}
\label{fig:sm-recovery}
\end{figure}


The first PIV Auto attempt ends with Result \code{0x02}, Detail \code{0x01}
(trust store empty). The ACU requests the SM diagnostic record, then reads the
Secure Messaging Certificate Signer object with PIVGETDATA, Object Identifier
\code{0x5FC122}, Tag Length zero, Offset zero, and Requested Length zero.
When the read succeeds, the reply contains the complete object, including
signer certificate tag \code{0x70} and any intermediate CVC in tag \code{0x7F21}.
The PD reads the credential through the contactless interface. If the object
is absent or the credential denies access, the PD reports the corresponding
PIVERROR. The ACU then selects another supported workflow or ends the
transaction. Recovery relies on the credential returning the object.

The ACU validates the signer certificate's chain to an authorized root,
validity period, revocation status, and PIV or PIV-I content-signing extended
key usage. The ACU checks any intermediate CVC's signature and linkage as
specified under PIVVCILOADTA. The anchor record contains the validated signer's
public key and the IIN derived from that certificate's subjectKeyIdentifier.
The ACU loads an anchor only after these checks succeed.
An IIN match alone provides no authorization to trust a certificate.

After the load succeeds, the ACU reapplies the PIVMODE configuration. The PD
starts a new attempt, establishes SM, and reports Result \code{0x05}, Detail
\code{0x01} (pairing code required). The ACU sends PIVXMITPAIRING and waits for
CARDSTATUSR with security status \code{0x01} (SM and VCI established).
The ACU then sends PIVSPE Mode \code{0x01}, with Auto Execute clear, to collect
and verify the PIN. After successful verification, the ACU reads the PIV
Authentication certificate (Object Identifier \code{0x5FC105}) and uses CRAUTH
with key \code{0x9A} and a fresh challenge. Certificate and challenge-response
validation remain the ACU's responsibility.

The ACU waits for a successful result at each step before issuing the next
command. An error ends this recovery sequence. These are commanded operations
after the automatic attempt has ended. The PD
returns each command's defined reply; the sequence generates no second PIV Auto
result. PIVMODE remains configured for the next credential presentation.
An ordinary PIVMODE replacement after pairing or PIN verification clears the state
needed for the commanded authentication. Alternatively, after the completed
pairing/PIN prerequisite failure, an eligible PIVMODE carrying the explicit Resume
TLV can preserve that context and start one new Auto attempt with the exact next
unused sequence counter. Subsection~\ref{sec:pivauto-resume} defines this distinct
operation; the illustrative resume bytes appear in the autonomous-template appendix.

If SM or VCI is unsupported, the ACU chooses a supported interface or
authentication method explicitly, or rejects the transaction. If a signature
or cryptogram check fails, the PD blocks credential operations and retains
the diagnostic record. The ACU investigates that failure before requesting
a new attempt. Trust recovery is appropriate for missing authorized trust
material; failed cryptographic checks are security failures.

This is a control-flow example. The preceding captured SM/VCI packet examples
cover pairing and PIN verification. The companion
\href{run:worked-examples/sm-vci/trust-recovery.md}{trust-recovery example}
records the command sequence and validation decisions.



\clearpage
\section{CRAUTH Worked Examples (Informative)}\label{app:crauth-examples}

This example sends a captured P-256 challenge to key \code{0x9E} and returns
the captured signature template. The first command fragment contains only
Format Marker \code{0xFF}. The second contains Algorithm \code{0x11},
Reference Identifier \code{0x9E}, and the \code{0x20}-byte challenge.
The complete request record is \code{0x0023} bytes. The PD removes the
marker and uses the algorithm, key reference, and challenge to construct
\code{GENERAL AUTHENTICATE}.

The credential data comes from the existing contactless card04 SM capture.
Credential SM is assumed established as in that capture; the example shows
the plaintext at the PD's SM boundary. The packets omit OSDP Secure Channel
for clarity. The two failure cases are modeled to show terminal error handling.
Full exchanges and source references are in
\href{run:worked-examples/crauth/examples.json}{the CRAUTH example file}.
The command and reply codes remain \code{0xA5} and \code{0x82}.

% Generated by scripts/render_crauth_examples.py; do not hand edit.
\begin{SimpleTable}{CRAUTH command and reply exchange}{3}
{|L{0.19\SimpleUsableWidth}|L{0.52\SimpleUsableWidth}|L{0.29\SimpleUsableWidth}|}
{\TableHeaderCell{Message} & \TableHeaderCell{Packet bytes (hex)} & \TableHeaderCell{Notes}}
ACU\newline \texttt{CRAUTH} & \texttt{53 01 0F 00 05 A5 23 00 00 00 01 00}\newline \texttt{FF 6B 63} & Total 0x0023; offset 0x0000; fragment 0x0001. Format marker only. \\\hline
PD\newline \texttt{ACK} & \texttt{53 81 08 00 05 40 39 35} & Fragment accepted. \\\hline
ACU\newline \texttt{CRAUTH} & \texttt{53 01 30 00 06 A5 23 00 01 00 22 00}\newline \texttt{11 9E 9B C1 80 EF 9A 61 6F 9F 05 8B}\newline \texttt{F4 60 44 C8 9A 30 1B BD 90 13 7A DF}\newline \texttt{12 F9 D9 A9 91 7B 09 4A EF A0 5C FF} & Total 0x0023; offset 0x0001; fragment 0x0022. \\\hline
PD\newline \texttt{ACK} & \texttt{53 81 08 00 06 40 6A 60} & Fragment accepted. \\\hline
ACU\newline \texttt{POLL} & \texttt{53 01 08 00 07 60 E9 55} & Poll for the result. \\\hline
PD\newline \texttt{CRAUTHR} & \texttt{53 81 36 00 07 82 4B 00 00 00 28 00}\newline \texttt{7C 49 82 47 30 45 02 21 00 83 6F 34}\newline \texttt{70 E4 CD 87 3C 6C A3 E3 B6 5A D3 6B}\newline \texttt{17 DF FC E2 7B 56 A8 25 EC 39 A9 DC}\newline \texttt{E3 D4 CB B0 6F 70} & Total 0x004B; offset 0x0000; fragment 0x0028. \\\hline
ACU\newline \texttt{POLL} & \texttt{53 01 08 00 05 60 8B 33} & Poll for the result. \\\hline
PD\newline \texttt{CRAUTHR} & \texttt{53 81 31 00 05 82 4B 00 28 00 23 00}\newline \texttt{95 02 20 61 1E FE 89 BF A8 13 B1 34}\newline \texttt{BD 11 1D 64 00 01 8B 50 71 1E 6C 9C}\newline \texttt{5D C8 22 DC 75 24 0B FA 60 D6 EE D9}\newline \texttt{FA} & Total 0x004B; offset 0x0028; fragment 0x0023. \\\hline
ACU\newline \texttt{POLL} & \texttt{53 01 08 00 06 60 D8 66} & Poll for the result. \\\hline
PD\newline \texttt{ACK} & \texttt{53 81 08 00 06 40 6A 60} & No completion pending. \\\hline
\end{SimpleTable}
\begin{SimpleTable}{Modeled failure completions}{3}
{|L{0.24\SimpleUsableWidth}|L{0.42\SimpleUsableWidth}|L{0.34\SimpleUsableWidth}|}
{\TableHeaderCell{Failure} & \TableHeaderCell{Packet bytes (hex)} & \TableHeaderCell{ACU action}}
Credential rejection & \texttt{53 81 0C 00 07 8A 23 10 69 82 2A D5} & PIVERROR 0x1023 (security conditions unmet), SW1/SW2 0x6982. End the request. \\\hline
Failure after a data fragment & \texttt{53 81 0A 00 05 8A 02 00 0E 51} & PIVERROR 0x0002 (smart-card timeout). Discard the partial response. \\\hline
\end{SimpleTable}


